\documentclass[10pt]{amsart}
\usepackage[a4paper,margin=23mm]{geometry}
\usepackage{amsmath,amssymb,amsthm,mathtools,hyperref,graphicx,xspace,letltxmacro,xpatch}
\usepackage[english]{babel}
\newtheorem{theo}{Theorem}[section]
\newtheorem{lemm}[theo]{Lemma}\newtheorem{prop}[theo]{Proposition}\newtheorem{coro}[theo]{Corollary}\newtheorem{conj}[theo]{Conjecture}
\theoremstyle{definition}\newtheorem{defi}[theo]{Definition}\newtheorem{rema}[theo]{Remark}\newtheorem{exam}[theo]{Example}\newtheorem{exem}[theo]{Example}\newtheorem{nota}[theo]{Notation}
\newcommand{\ie}{i.e.}\newcommand{\eg}{e.g.}\newcommand{\resp}{resp.}\newcommand{\skpt}{\par\smallskip\noindent}
\emergencystretch=3em\pagestyle{plain}

\usepackage{mathrsfs}

\DeclareMathOperator{\supp}{supp}
\DeclareMathOperator{\diver}{div}
\DeclareMathOperator{\Osc}{Osc}

\newcommand{\p}{\partial}
\newcommand{\e}{\varepsilon}
\newcommand{\mek}{{\mathbf 1}}
\newcommand{\we}{\mathfrak{w}}
\newcommand{\ph}{\varphi}
\newcommand{\R}{{\mathbb R}}
\newcommand{\N}{{\mathbb N}}
\newcommand{\Z}{{\mathbb Z}}
\newcommand{\IP}{{\mathbb P}}
\newcommand{\pP}{{\mathbb P}}
\newcommand{\I}{{\mathbb I}}
\newcommand{\E}{{\mathbb E}}
\newcommand{\T}{{\mathbb T}}
\newcommand{\CC}{{\mathcal C}}
\newcommand{\DD}{{\mathcal D}}
\newcommand{\EE}{{\mathcal E}}
\newcommand{\FF}{{\mathcal F}}
\newcommand{\MM}{{\mathcal M}}
\newcommand{\PP}{{\mathcal P}}
\newcommand{\RR}{{\mathcal R}}
\newcommand{\VV}{{\mathcal V}}
\newcommand{\vk}{{\varkappa}}
\newcommand{\la}{\lambda}
\newcommand{\ty}{\infty}
\newcommand{\om}{\omega}
\newcommand{\de}{\delta}
\newcommand{\wwww}{{\mathfrak w}}
\newcommand{\mm}{{\mathfrak m}}
\newcommand{\elll}{{\mathcal L}}
\newcommand{\pp}{{{\mathsf P}_N}}
\newcommand{\qq}{{{\mathsf Q}_N}}
\newcommand{\PPPP}{{\mathfrak P}}
\newcommand{\BBBB}{{\mathfrak B}}
\newcommand{\BBBBB}{{\mathcal B}}
\newcommand{\SSS}{{\mathscr S}}

\newcommand{\dt}{\dot}

\newcommand{\lag}{\langle}
\newcommand{\rag}{\rangle}

\newcommand{\dd}{\mathrm{d}}

\newtheorem*{mtheorem}{Main Theorem}

\newenvironment{step}{\begin{proof}}{\let\qed\relax \end{proof}}



\graphicspath{{./figures/}}

\newcommand*{\mk}{\mkern -1mu}
\newcommand*{\Mk}{\mkern -2mu}
\newcommand*{\mK}{\mkern 1mu}
\newcommand*{\MK}{\mkern 2mu}

\hypersetup{urlcolor=purple, linkcolor=blue, citecolor=red}

\let\oldexists\exists
\renewcommand*{\exists}{\mathrel{\oldexists}}
\let\oldforall\forall
\renewcommand*{\forall}{\mathrel{\oldforall}}

\newenvironment{proofc}[1][\proofname]{\begin{proof}[{#1}]}{\end{proof}}

\providecommand{\pointir}{.\enspace---\enspace}
\title{Weighted multiplicative ergodicity of Feynman-Kac semigroups for stochastic Navier-Stokes}
\author{Nersesyan, Vahagn}
\date{doi:10.5802/ahl.23; CC BY 4.0}
\begin{document}
\maketitle
\noindent\textit{Editorial note.} Only selected author theorem counted; complete original mathematical body and bibliography preserved. Native vector figures retained; publisher date/editor frontmatter omitted from rendering preamble and proofc class environment supplied where needed.\par\medskip
\noindent\textit{Editorial note.} Native pointir heading separator supplied as period and em dash in rendering wrapper.\par\medskip
\noindent\textit{Allowed standard prerequisites (editorial).}\begin{itemize}
\item Standard functional analysis, measure theory and the author-cited classical theorems. Standard prior results stay cited; complete source-local proofs and definitions retained.
\end{itemize}\medskip

\section{Introduction}\label{S:0}

We study the large deviations principle (LDP) for the 2D Navier--Stokes system for incompressible fluids:
\begin{equation}\label{0.1}
\p_t u+\langle u,\nabla \rangle u-\nu\Delta u+\nabla p=f(t,x), \quad\diver u=0,
\end{equation}
where $\nu>0$ is the viscosity of the fluid, $u = (u_1(t,x), u_2(t,x))$ and~$p = p(t, x)$ are unknown velocity field and pressure, $f$ is an external (random) force, and $\lag u, \nabla \rag=u_1\partial_1+u_2\partial_2$. Throughout this paper, we assume that the space variable~$x = (x_1, x_2)$ belongs\,\footnote{The periodic boundary conditions are chosen to simplify the presentation. Similar results can be established in the case of a bounded domain with smooth boundary and Dirichlet boundary conditions.} to the standard torus~$\T^2=\R^2/2\pi\Z^2$. The problem is considered in the space of divergence-free vector fields with zero mean~value
\begin{equation}\label{space1}
H= \left\{u\in L^2(\T^2,\R^2):\,\, \diver u=0 \,\, \text{in $\T^2$,} \,\,\,\int_{\T^2} u(x) \dd x=0 \right\}
\end{equation}
endowed with the $L^2$-norm $\|{\,\cdot\,}\|$. We assume that the force is of the~form
\[
f(t,x)=h(x)+ \eta(t,x),
\]
where $ h\in H^1:= H^1(\T^2, \R^2)\cap H$ is a given function and~$\eta$ is a white-in-time~noise
\begin{equation}\label{0.2}
\eta(t,x)= \p_t W(t,x),\quad W(t,x)=\sum_{j=1}^\ty b_j\beta_j(t) e_j(x).
\end{equation}
Here $\{b_j\}$ is sequence of real numbers such that
\begin{equation}
{\BBBB_{1}}=\sum_{j=1}^\ty \alpha_j b_j^2<\ty, \label{0.3}
\end{equation}
$\{\beta_j\} $ is a sequence of independent standard Brownian motions defined on a filtered probability space\,\footnote{We assume that this space satisfies the usual conditions (see~\cite[Definition~2.25]{KS1991}).} $(\Omega,\FF, \{\FF_t\}, \pP)$, and $\{e_j\}$ is an orthonormal basis in $H$ consisting of the eigenfunctions of the Stokes operator~$L=- \Delta $ with eigenvalues $\{\alpha_j\}$.~As usual, projecting~\eqref{0.1} to $H$, we eliminate the pressure and obtain an evolution equation for the velocity field\,\footnote{To simplify the notation, we shall assume that $\nu = 1$.} (e.g., see~\cite[Chapter~1, Section~6]{lions1969}):
\begin{equation}\label{0.4}
\dot u+B(u)+L u =h(x)+ \eta(t,x),
\end{equation}
where $B(u)=\Pi(\langle u,\nabla \rangle u)$ and $\Pi$ is the orthogonal projection onto $H$ in~$L^2$. This system is supplemented with the initial condition
\begin{equation}\label{E:init}
u(0)=u_0.
\end{equation}
Under these assumptions, problem~\eqref{0.4}, \eqref{E:init} admits a unique solution and defines a Markov family~$(u_t,\IP_u)$ parametrised by the initial condition~$u=u_0 \in H$. The ergodic properties of this family are now well understood. In particular, it is known that~$(u_t,\IP_u)$ admits a unique stationary measure, which is exponentially mixing, provided that sufficiently many coefficients $b_j$ are non-zero (see the papers~\cite{FM-1995,KS-cmp2000,EMS-2001, KS-jmpa2002, BKL-2002,HM-2006, odasso-2008} and the book~\cite{KS-book}). A central limit theorem~(CLT) for problem~\eqref{0.4},~\eqref{E:init} is established in~\cite{kuksin-jmpa, shirikyan-ptrf}. The LDP proved in the present paper is a natural extension of the CLT. Indeed, while the CLT describes the probability of small deviations of a time average of a functional from its mean value, the LDP quantifies the probability of large deviations.

%\smallskip

Before formulating the main result of this paper, let us introduce some notation and definitions.~We denote by $\PP(H)$ the space of Borel probability measures on~$H$ endowed with the to\-pology of weak convergence.~Given a measure~$\nu\in\PP(H)$, we set $
\IP_\nu(\Gamma)=\int_H\IP_u(\Gamma)\nu(\dd u)$ for any~$\Gamma\in \FF$ and consider the following family of \emph{occupation~measures}
\begin{equation}\label{0.5}
\zeta_t=\frac1t\int_{0}^{t}\delta_{u_s} \dd s,\quad t>0
\end{equation}
defined on the probability space $(\Omega,\FF,\IP_\nu)$. Here~$\de_u$ is the Dirac measure concentrated at $u\in H$. We shall say that a mapping $I : \PP(H) \to [0, +\ty]$ is a \emph{good rate function} if the level set~$\{\sigma\in\PP(H) : I(\sigma) \le \alpha \}$ is compact for any~$\alpha \ge0$.~A good rate function~$I$ is nontrivial if its effective domain $D_I = \{\sigma \in \PP(H) : I(\sigma) < \ty\}$ is not a singleton. For any numbers~$\varkappa>0$ and~$M>0$, we denote
\[
\Lambda(\varkappa,M)=\left\{\nu\in\PP(H):\int_He^{\varkappa\|v\|^2}\nu(\dd v)\le M \right\}.
\]

\begin{mtheorem}
Assume that~\eqref{0.3} holds and $b_j>0$ for all $j\ge1$.~Then for any numbers $\vk>0$ and $M>0$, the family~$\{\zeta_t,t>0\}$ satisfies an LDP, uniformly~with respect to $\nu\in\Lambda(\varkappa,M)$, with a non-trivial good rate function $I: \PP(H)\to [0,+\ty]$ not depending on~$\vk$ and $M$. More precisely, the following two bounds~hold.

Upper bound\pointir For any closed subset~$F\subset\PP(H)$, we have
\[
\limsup_{t\to\infty} \frac1t\log \sup_{\nu\in\Lambda} \pP_\nu\{\zeta_t\in F\}\le -\inf_{\sigma\in F} I(\sigma).
\]

Lower bound\pointir For any open subset~$G\subset\PP(H)$, we have
\[
\liminf_{t\to\infty} \frac1t\log \inf_{\nu\in\Lambda}\pP_\nu\{\zeta_t\in G\}\ge -\inf_{\sigma\in G}I(\sigma).
\]

Furthermore, $I$ is given by
\begin{equation}\label{E:II}
I(\sigma)=\sup_{V\in C_b(H)} \left(\int_H V(u)\sigma(\dd u)-Q(V) \right),
\quad \sigma\in\PP(H),
\end{equation}
where~$Q:C_b(H)\to\R$ is a~$1$-Lipschitz convex function such that~$Q(C)=C$ for any~$C\in\R$.
\end{mtheorem}

This type of large deviations results have been first established by Donsker and Varadhan~\cite{DV-1975} and later generalised by many others (see the books~\cite{FW1984,DS1989,ADOZ00} and the references therein). There are only a few works studying the large deviations behaviour of solutions of randomly forced PDEs as time goes to infinity. The case of the stochastic Burgers and Navier--Stokes equations is first studied in~\cite{gourcy-2007b, gourcy-2007a}. In these papers, the random perturbation is of the form~\eqref{0.2} with the following restriction on the coefficients
\begin{equation}\label{Gougz}
cj^{-\alpha}\le b_j\le Cj^{-\frac12-\e}, \quad \frac12<\alpha<1,\,\, \e\in \left(0, \alpha-\frac12\right].
\end{equation}
Notice that the lower bound does not allow the sequence $\{b_j\}$ to converge to zero sufficiently fast, so the external force $f$ is \emph{irregular} with respect to the space variable. This is not very natural from the physical point of view. The proof is based on a general sufficient condition established in~\cite{wu-2001}, and essentially uses the \emph{strong Feller property}. The main novelty of our Main Theorem is that it proves an LDP without any lower bound on $\{b_j\}$ (so, in particular, we do not have a strong Feller property).

%\smallskip
We use an approach introduced in the papers~\cite{JNPS-2012, JNPS-2014}, where an LDP is established for a family of dissipative PDEs perturbed by a \emph{random kick force}. The proofs of these papers are based on a Kifer type criterion for LDP and a study of the large-time behaviour of generalised Markov semigroups. These results have been later extended in~\cite{MN-2015} to the case of the stochastic damped nonlinear wave equation driven by a \emph{spatially regular white noise}. The main result of that paper is an LDP of local type.~In the case of the Navier--Stokes system~\eqref{0.4}, although we follow a similar scheme, there are important~differences in all the steps of the argument, coming from both the continuous-time nature of the system and the globalness of the LDP. Here we study the large-time asymptotics of the Feynman--Kac semigroup without any restriction on the smallness of the potential.~One of the most important~difficulties arises in the proof of the uniform Feller property.~To establish this,~we construct coupling processes using \emph{a~new~two parameter auxiliary equation} (see~\eqref{auxi}) which allows to have an appropriate Foia\c{s}--Prodi~estimate for~the~trajectories and a rapid exponential stabilisation for finite-dimensional~projections.

%\smallskip

Let us also mention that~the multiplicative ergodic theorem we obtain for system~\eqref{0.4} is of more general~form and works for a larger class of functionals and initial measures (see~Theorem~\ref{T:1.1}).

It is a challenging \emph{open problem} whether an LDP still holds for~\eqref{0.4}, \eqref{E:init} when the driving noise is \emph{highly degenerate} (i.e., only a finite number of $b_j$ are non-zero in~\eqref{0.2}). For the Navier--Stokes system in this degenerate situation, exponential mixing is established in~\cite{HM-2006} for white-in-time noise and in~\cite{KNS-2018} for a bounded noise satisfying some decomposability and observability hypotheses. Using these results and literally repeating the arguments of the proof of Theorem~5.4 in~\cite{MN-2017}, one can prove a level-1 LDP of local type.

%\smallskip
The paper is organised as follows.~In Section~\ref{S:1}, we state a multiplicative~ergodic theorem for the Navier--Stokes system and combine it with Kifer's criterion for non-compact spaces to prove the Main Theorem. In Sections~\ref{S:2} and~\ref{S:3}, we check the conditions of an abstract result on large-time behaviour of generalised Markov semigroups. Section~\ref{S:3Sd} is devoted to the proof of the multiplicative~ergodicity.~In the Appendix, we prove various a priori estimates for the~solutions and~recall the statement of the above-mentioned result for generalised Markov~semigroups.

\subsection*{Notation}

We shall use the following standard notation.
\begin{itemize}
%\smallskip\noindent
\item
$H$ is the space defined by~\eqref{space1}, $B_H(a,R)$ is the closed ball in~$H$ of radius~$R$ centred at~$a$. When~$a=0$, we write $B_H(R)$.

%\smallskip\noindent
\item
$H^1= H^1(\T^2, \R^2)\cap H$, where $H^1(\T^2, \R^2)$ is the space of vector functions $u = (u_1,u_2)$ with components in the usual Sobolev space of order $1$ on $\T^2$.

%\smallskip\noindent
\item
$L^\infty(H)$ is the space of bounded Borel-measurable functions $f:H\to\R$ endowed~with the norm $\|f\|_\infty=\sup_{u\in H}|f(u)|$.

%\smallskip\noindent
\item
$C_b(H)$ is the space of continuous functions~$f\in L^\infty(H)$.

%\smallskip\noindent
\item
$L_b(H)$ is the space of functions $f\in C_b(H)$ for which the following norm is finite
\[
\|\psi\|_{L}=\|\psi\|_\infty+\sup_{u\neq v} \frac{ {|\psi(u)-\psi(v)|}}{ \|u-v\| }.
\]

%\smallskip\noindent
\item
$\VV$ is the space of functions~$V:H\to \R$ for which there is an integer~$N \ge 1$ and a function~$F \in L_b(H_N)$ such that
\begin{equation}\label{repres}
V(u)=F(\pp u),\quad u\in H.
\end{equation}
Here $\pp$ is the orthogonal projection in~$H$ onto the space
\begin{equation}\label{ddftgr}
H_N=\text{span}\{e_1,\ldots, e_N\}
\end{equation}
and $\{e_j\}$ is the orthonormal basis entering~\eqref{0.2}.

%\smallskip\noindent
\item
For a given Borel-measurable function ${\wwww}:H\to[1,+\infty]$, we denote by $C_\wwww(H)$~(respectively, $L_\wwww^\infty(H)$) the space of continuous (Borel-measurable) functions $f:H\to\R$ such that
\[
\|f\|_{L_\wwww^\infty}=\sup_{u\in H}\frac{|f(u)|}{\wwww(u)}<\ty.
\]

%\smallskip\noindent
\item
$\MM_+(H)$ is the set of non-negative finite Borel measures on $H$ endowed~with the topology of the weak convergence. For $\mu\in \MM_+(H)$ and~$f\in C_b(H)$, we denote~$\lag f,\mu\rag=\int_Hf(u)\mu(\dd u).$

%\smallskip\noindent
\item
$\PP(H)$ is the set of Borel probability measures on~$H$, and $\PP_\wwww(H)$ is the set of measures $\mu\in\PP(H)$ such that $\langle \wwww,\mu\rangle<\infty$.
\end{itemize}


\section{Proof of the Main Theorem}\label{S:1}

In this section, we state a multiplicative ergodic theorem for the Navier--Stokes system~\eqref{0.4} and apply it to prove the Main Theorem. Let us start by introducing the following two \emph{weight functions}
\begin{align}
\mm_\vk(u)&=\exp(\vk \|u\|^2), \label{1.2}
\\
\wwww_m(u)&=1+\|u\|^{2m},\quad u\in H \label{1.1}
\end{align}
for positive numbers $\varkappa$ and $m$. To avoid triple subscripts, we shall write~$C_{\mm}(H)$ and $\PP_{\mm}(H)$ instead of~$C_{\mm_\vk}(H)$ and~$\PP_{\mm_\vk}(H)$. Recall that the \emph{Feynman--Kac semigroup} associated with the family~$(u_t,\IP_u)$ is defined by
\[
\PPPP_t^V f(u)=\E_u \left\{ \Xi^V_t f(u_t)\right\},
\]
where
\begin{equation}\label{ARM}
\Xi^V_t = \exp\left(\int_{0}^tV(u_s) \dd s\right).
\end{equation}
From estimate~\eqref{5.2} it follows that, for sufficiently small $\vk$ and any $V\in C_b(H)$, the application~$\PPPP_t^V$ maps $C_{\mm}(H)$ into itself.~Let $\PPPP_t^{V*}:\MM_+(H)\to \MM_+(H)$~be its dual. Then a measure~$\mu\in\PP(H)$ and a function $h\in C_\mm(H)$ are \emph{eigenvectors} corresponding to an eigenvalue $\la>0$ if
\[
\PPPP^{V*}_t\mu=\la^t \mu, \qquad \PPPP^V_t h=\la^t h\quad \text{for any~$t>0$}.
\]
We have the following result.


\begin{theo}\label{T:1.1}
Under the conditions of the Main Theorem, for any $V\in\VV$, there are numbers $m=m(V)\ge1$ and~$\gamma_0=\gamma_0(\BBBB_{0})>0$, where $\BBBB_{0}=\sum_{j\ge1}b_j^2$, such that the following assertions hold for any $\vk\in (0,\gamma_0).$

Existence and uniqueness\pointir There is a unique pair $(\mu_V,h_V)\!\in\! \PP_\mm(H)\times C_\wwww(H)$ of eigenvectors corresponding to an eigenvalue $\la_V>0$ normalised by the condition $\langle h_V,\mu_V\rangle=1$.

Convergence\pointir For any $f\in C_{\mm}(H)$, $\nu\in \PP(H)$, and $R>0$, we have
\begin{align}
\lambda_V^{-t}\PPPP_t^V f&\to\lag f,\mu_V\rag h_V \quad\text{in~$C_b(B_H(R))\cap L^1(H,\mu_V)$ as~$t\to\infty$}, \label{1.5}
\\
\lambda^{-t}_V\PPPP_t^{V*}\nu&\to\lag h_V,\nu\rag\mu_V
\quad\text{in~$\MM_+(H)$ as~$t\to\infty$}. \label{1.6}
\end{align}
Moreover, for any $M>0$ and $\vk'\in (\vk,\gamma_0)$, the convergence
\begin{equation}\label{1.7}
\lambda_V^{-t}\,\E_\nu\left\{f(u_t)\exp\bigg(\int_0^tV(u_s)\dd s\bigg)\right\}\to\langle f,\mu_V\rangle\,\langle h_V,\nu\rangle\quad \text{as } t\to\infty
\end{equation}
holds uniformly in $\nu\in\Lambda(\vk',M)$.
\end{theo}


This theorem is established in Section~\ref{S:3Sd}.~Here we combine it with some~arguments from~\cite{JNPS-2014, MN-2015}, to prove the Main Theorem.

\begin{proofc}[Proof of the Main Theorem]

\begin{step}[Step~1: Reduction]
It suffices to prove the Main Theorem for small $\vk$, so we shall assume that $\vk\in (0, \gamma_0)$.
Let us take any~$M>0$ and endow the set
\[
\Theta= \R^*_+\times\Lambda(\vk,M)
\]
with an order relation $\prec$ defined by $(t_1,\nu_1)\prec (t_2,\nu_2)$ if and only if~$t_1\le t_2$.~Then a family~$\{x_\theta\in \R,\theta\in\Theta\}$ converges if and only if it converges uniformly~with~respect to~$\nu\in\Lambda(\vk,M)$~as~$t\to\ty$.~Assume that the following three properties~hold.
\begin{enumerate}
\item \label{prop1} For any $V\in C_b(H)$ and~$\nu\in \Lambda(\vk,M)$, the following limit exists
\[
Q(V)=\lim_{t\to\ty} \frac{1}{t}
\log \E_\nu \exp \biggl(\,\int_0^tV(u_s)\dd s\biggr).
\]
Moreover, it does not depend and is uniform in $\nu\in \Lambda(\vk,M)$.

\item \label{prop2} There is a vector space $\VV\subset C_b(H)$ such that its restriction to any compact set $K\subset H$ is dense in $C(K)$, and for any~$V\in \VV $, there is a unique~$\sigma_V\in \PP(H)$ satisfying the relation
\begin{equation}\label{equilhdd}
Q(V)= \lag V, \sigma_V\rag-I(\sigma_V),
\end{equation}
where $I(\sigma)$ is the Legendre transform of~$Q$ given by~\eqref{E:II}.

\item \label{prop3} There is a function $\varPhi : H \to [0,+\ty]$ with compact level sets $ \{u \in H : \varPhi(u)\le \alpha\}$ for any $\alpha\ge0$ such that
\begin{equation}\label{1.8}
\E_\nu\exp \biggl(\,\int_0^t\varPhi(u_s)\dd s\biggr)\le Ce^{c t},\quad \nu\in\Lambda(\vk,M),\, t>0
\end{equation}
for some positive constants $C$ and $c$.
\end{enumerate}
For any $\theta=(t,\nu)\in \Theta$, let us set $r_\theta:=t$ and~$\zeta_\theta:=\zeta_t$, where $\zeta_t$ is the random probability measure given by~\eqref{0.5} defined on the probability space~$(\Omega_\theta,\FF_\theta,\IP_\theta):=(\Omega,\FF,\IP_\nu)$. The definition of the relation~$\prec$ and Properties~\eqref{prop1}--\eqref{prop3} imply that the family~$\{\zeta_\theta\}$ satisfies the conditions of the Kifer type criterion established in Theorem~3.3 in~\cite{JNPS-2014}. Hence~\eqref{E:II} defines a good rate function $I$ and for any closed set $F\subset \PP(H)$ and open set $G\subset \PP(H)$, we have
\begin{align*}
\limsup_{\theta\in\Theta} \frac{1}{r_\theta}\log\pP_\theta\{\zeta_\theta\in F\}&\le -\inf_{\sigma\in F} I (\sigma),\\
\liminf_{\theta\in\Theta} \frac1{r_\theta}\log\pP_\theta\{\zeta_\theta\in G\}&\ge -\inf_{\sigma\in G} I (\sigma).
\end{align*}
These two inequalities imply the upper and lower bounds in the Main Theorem, since we have the following equalities
\begin{align*}
\limsup_{\theta\in\Theta} \frac{1}{r_\theta}\log\pP_\theta\{\zeta_\theta\in F\}&=\limsup_{t\to\infty} \frac1t\log \sup_{\nu\in\Lambda} \pP_\nu\{\zeta_t\in F\},\\
\liminf_{\theta\in\Theta} \frac1{r_\theta}\log\pP_\theta\{\zeta_\theta\in G\}&=\liminf_{t\to\infty} \frac1t\log \inf_{\nu\in\Lambda}\pP_\nu\{\zeta_t\in G\}.
\end{align*}
Now we turn to the proofs of Properties~\eqref{prop1}--\eqref{prop3}.
\end{step}

\begin{proof}[Step~2: Proof of Properties~\eqref{prop1}--\eqref{prop3}]
Property~\eqref{prop3} is the easiest one. It is verified for $\varPhi(u)=\varkappa \|u\|_1^2$ if we choose~$\varkappa\in (0,\gamma_0)$. Indeed,~\eqref{1.8} follows from inequality\,\footnote{We shall see in the proof of Theorem~\ref{T:1.1}, that $\gamma_0$ is the number in Lemma~\ref{L:5.1}.}~\eqref{5.1}, and $\varPhi$~has compact level sets, since it is continuous from~$H^1$ to $\R$ and the embedding $H^1\subset H$ is compact.

Properties~\eqref{prop1} and~\eqref{prop2} are proved using the same methods as in the case of the discrete-time model considered in~\cite{JNPS-2014}. The restriction of $\VV$ to any compact set~$K\subset H$ is dense in $C(K)$. Taking $f=\mek $ in~\eqref{1.7}, we get Property~1 for any~$V\in \VV$ with~$Q(V)=\log \la_V$. In the case of an arbitrary~$V\in C_b(H)$, this property is established by using a \emph{buc-approximating} sequence~$V_n\in\VV$ of $V$ (i.e.,~$\sup_{n\ge1}\|V_n\|_\ty \!< \!\ty$ and $\|V_n-V\|_{L^\ty(K)}\to 0$ as $ n\to\ty$ for any compact $K$ in $ H$) and the exponential tightness of the family~$\{\zeta_\theta\}$ (which holds by Property~3). The reader is referred to Section~5.6 of~\cite{JNPS-2014} for the details.

%\smallskip
To prove Property~\eqref{prop2}, for any $V\in \VV$ and $F\in C_b(H)$, we consider the following auxiliary Markov semigroup
\[
\SSS_t^{V,F}g (u)=\lambda_V^{-t}h_V^{-1}\PPPP_t^{V+F}(h_V g)(u), \quad g\in C_b(H), \, t\ge0.
\]
By Property~\eqref{prop1}, the following limit exists
\[
Q^V(F):=\lim_{t\to\ty} \frac{1}{t}\log(\SSS_t^{V,F}{\mek})(u).
\]
Let $I^V: \MM(H)\to [0,+\infty] $ be the Legendre transform of $Q^V$.
The arguments of Section~5.7 in~\cite{JNPS-2014} show that $\sigma\in \PP(H)$ satisfies~\eqref{equilhdd} if and only if $I^V(\sigma)=0$. On the other hand, by Proposition~1.3 in~\cite{MN-2015} (whose proof is the same in our case), the measure~$\sigma_V=h_V\mu_V$ is the unique zero of~$I^V$.

%\smallskip

It remains to show that the good rate function $I$ is non-trivial. Assume, by contradiction, that~$D_I$ is a singleton.~Then $I(\mu)=0$ and $I(\sigma)=+\ty$ for $\sigma\in\PP(H)\setminus\{\mu\}$, where $\mu$ is the stationary measure of $(u_t,\IP_u)$. On the other hand, as the Legendre transform is its own inverse, we derive from~\eqref{E:II} that
\[
Q(V)= \sup_{\sigma\in\PP(H)} \bigl(\lag V, \sigma\rag-I(\sigma)\bigr) \quad \text{for $V\in C_b(H)$}.
\]
This implies that $Q(V)=\lag V,\mu\rag$ for any $V\in C_b(H)$. Let us take any non-constant $V\in \VV$ such that $\lag V,\mu\rag=0$. Then $Q(V)=0$, and from limit~\eqref{1.5} with $f={\mathbf1}$ and~$\nu=\mu$ we get~$\la_V=e^{Q(V)}=1$ and
\begin{equation}\label{e52}
\sup_{t\ge0}\E_\mu \exp\left(\int_0^t V(u_s)\dd s\right)<\infty.
\end{equation}
Combining the latter with the central limit theorem (see~\cite[Proposition~4.1.4]{KS-book}), we get $V= {\mathbf0}$.~This contradicts the assumption that $V$ is non-constant and completes the proof of the Main Theorem.
\end{proof}
\let\qed\relax
\end{proofc}


\section{Checking conditions of Theorem~\ref{T:5.3}}\label{S:2}

Theorem~\ref{T:1.1} is proved by applying a convergence result for generalised Markov semigroups obtained in~\cite{JNPS-2014, MN-2015} and restated here as Theorem~\ref{T:5.3}. In this and next sections, we show that the conditions of that theorem are satisfied for the generalised Markov family of transition kernels defined by
\[
P_t^V(u,\Gamma)=(\PPPP_t^{V*} \de_u) (\Gamma),\quad \Gamma\in \BBBBB(H), \, u\in H,
\]
if we take $X=H$, $X_R= B_{H^1}(R)$, and~$\wwww= \wwww_m$ with sufficiently large $m\ge1$.

\subsection{Growth estimates}

Estimate~\eqref{5.5} implies that the measure~$P_{t}^V(u,{\cdot\,})$ is concentrated on the space $H^1=\bigcup_{R=1}^\ty X_R=X_\ty$ for any $V\in C_b(H)$, $t>0$, and $u\in H$. The boundedness of~$V$ implies that $\sup_{t\in[0,1]}\|\PPPP_t^V{\mek}\|_\ty<\ty$. So the following proposition gives the growth condition in Theorem~\ref{T:5.3}.

\begin{prop}\label{P:2.4}
For any $V\in C_b(H)$, there are positive numbers $m$ and $R_0$ such that
\begin{align}
&\sup_{t\ge0}
\frac{\|\PPPP_t^V\wwww\|_{L_{\wwww}^\infty}}{\|\PPPP_t^V{\mek}\|_{R_0}}<\infty,\label{a5.8}
\end{align}
where $\wwww=\wwww_m$ and~$\|{\,\cdot\,}\|_{R_0}$ is the $L^\ty$ norm on $X_{R_0}$.
\end{prop}

\begin{proof}
Replacing $V$ by $V-\inf_HV$, we can assume that~$V$ is non-negative.

\begin{step}[Step 1]
Let us show that there are integers $m, R_0\ge1$ such that
\begin{equation}\label{6.0015}
\sup_{t\ge0} \frac{\|\PPPP_{t}^V{\mek}\|_{L_\wwww^\infty}}{\|\PPPP_{t}^V{\mek}\|_{R_0}}<\infty.
\end{equation}
Indeed, let $\tau(R)$ be the first hitting time of the set $X_R$ defined by~\eqref{E:tau}, and let $m$ and $R_0$~be~the integers in Proposition~\ref{P:2.1} for $\gamma=\|V\|_\ty$. Then for~any~$u\in H$, we have
\begin{equation}
\PPPP_t^V{\mek}(u)=\E_u\Xi^V_t =\E_u \bigl\{\I_{G_t}\Xi^V_t \bigr\}+\E_u\bigl\{\I_{G_t^c}\Xi^V_t\bigr\}=:I_1+I_2,\label{6.16}
\end{equation}
where $\Xi^V_t$ is given by~\eqref{ARM} and $G_t=\{\tau(R_0)> t\}$. As $V$ is non-negative, we have $\PPPP_t^V{\mek}(u)\ge1$. This and~\eqref{8.1} imply that
\begin{equation}\label{sdqsd}
I_1\le \E_u \Xi^V_{\tau(R_0)} \le \E_u\exp\bigl(\gamma\tau(R_0)\bigr)\le C\,\wwww(u)
\le C\,\wwww(u)\,\|\PPPP_t^V{\mek}\|_{R_0}.
\end{equation}
By the strong Markov property and~\eqref{8.1},
\begin{align}
I_2&\le \E_u\bigl\{\I_{G_t}\Xi^V_{\tau(R_0)}\E_{u(\tau(R_0))}\Xi^V_t\bigr\}
\nonumber\\&\le\E_u \{e^{\gamma\tau(R_0)} \}\,\|\PPPP_t^V{\mek}\|_{R_0} \le C\,\wwww(u)\,\|\PPPP_t^V{\mek}\|_{R_0}.\label{qqdfdz}
\end{align}
Inequalities~\eqref{6.16}--\eqref{qqdfdz} imply~\eqref{6.0015}.
\end{step}
%\smallskip

\begin{step}[Step 2]
It suffices to prove~\eqref{a5.8} for integer times $k\ge1$:
\begin{equation}
\sup_{k\ge0}
\frac{\|\PPPP_k^V\wwww\|_{L_{\wwww}^\infty}}{\|\PPPP_k^V{\mek}\|_{R_0}}<\infty.\label{a5.8b}
\end{equation}
Indeed, the semigroup property and the fact that $V$ is non-negative and bounded imply that
\begin{align*}
\|\PPPP^V_t\wwww\|_{L_\wwww^\infty} &=\|\PPPP^V_{t-[t]}(\PPPP^V_{[t]}\wwww)\|_{L_\wwww^\infty} \le C_0\|\PPPP^V_{[t] }\wwww\|_{L_\wwww^\infty},\\
\|\PPPP_t^V{\mek}\|_{R_0}&\ge\|\PPPP_{[t]}^V{\mek}\|_{R_0},
\end{align*}
where $[t]$ is the integer part of $t$ and $ C_0:=\sup_{s\in[0,1]}\|\PPPP^V_s\wwww\|_{L_\wwww^\infty}.$ By~\eqref{5.3}, we~have
\[
C_0 \le e^{\gamma} \sup_{s\in[0,1]}\|\PPPP_s\wwww\|_{L_\wwww^\infty}<\ty,
\]
where $\PPPP_t=\PPPP_t^0$ is the Markov operator associated with~\eqref{0.4}.
\end{step}

\begin{proof}[Step 3]
To prove~\eqref{a5.8b}, we use the Markov property and~\eqref{5.3}:
\begin{align*}
\PPPP_{k}^V\wwww(u) &\le e^{\gamma} \E_u \left\{\Xi^V_{k-1} \wwww(u_{k}) \right\}\\&= e^{\gamma} \E_u \left\{\Xi^V_{k-1} \E_{u_{k-1}} \wwww(u_1) \right\} \\&\le e^{\gamma} \E_u \left\{\Xi^V_{k-1} [e^{-m\alpha_1 }\we(u_{k-1})+C_1] \right\} \\&\le q \PPPP_{k-1}^V\wwww(u)+ e^{\gamma} C_1 \PPPP_{k-1}^V{\mek}(u),
\end{align*}
where we choose $m > \gamma/\alpha_1 $, so that $q:=e^{\gamma-m\alpha_1}<1$. Iterating this inequality and using the fact that the sequence $\{\|\PPPP_{k}^V{\mek}\|_{R_0}\}$ is a non-decreasing in $k$, we~obtain
\begin{align*}
\PPPP_{k}^V\wwww(u) \le q^k \wwww(u)+ (1-q)^{-1}e^{\gamma} C_1 \PPPP_{k}^V{\mek}(u).
\end{align*}
This and~\eqref{6.0015} imply~\eqref{a5.8b}.
\end{proof}
\let\qed\relax
\end{proof}

We shall also need the following growth estimates with two other weights.

\begin{prop}\label{C:2.5}
Let $V\in C_b(H)$ and let $R_0$ and $\gamma_0$ be the numbers in Proposition~\ref{P:2.4} and Lemma~\ref{L:5.1}, respectively. Then for any~$\vk\in(0,\gamma_0)$, we have
\begin{align}
&\sup_{t\ge0} \frac{\|\PPPP_t^V\mm\|_{L_\mm^\infty}}{\|\PPPP_t^V{\mek}\|_{R_0}}<\infty, \label{dqdlkk}
\\
&\sup_{t\ge1} \frac{\|\PPPP_t^V F \|_{L_{\mm}^\infty}}{\|\PPPP_t^V{\mek}\|_{R_0}}<\infty, \label{estimE}
\end{align}
where $\mm=\mm_\vk$ and $F(u)= \|u\|_1^2$.
\end{prop}

\begin{proofc}

\begin{step}[Step~1: Proof of~\eqref{dqdlkk}]
As in the previous proof, we can assume that~$V$ is non-negative and $t=k$ is integer. We take any $A>0$ and write
\begin{align}\label{aaerfcvfd}
\PPPP_k^V\mm (u)&= \E_u \left\{ \I_{\{\|u_k\|^2\le A\}} \Xi_k^V\mm(u_k) \right\}+ \E_u \left\{ \I_{\{\|u_k\|^2> A\}} \Xi_k^V\mm(u_k) \right\}\nonumber\\& =:I_k+J_k.
\end{align}
By~\eqref{6.0015}, we have
\[
\|\PPPP_k^V{\mek}\|_{L_\mm^\infty} \le C_2 \|\PPPP_{k}^V{\mek}\|_{R_0},
\]
hence
\begin{equation}\label{azzezfd}
\|I_k\|_{L_\mm^\infty} \le e^{\vk A} \|\PPPP_k^V{\mek}\|_{L_\mm^\infty} \le C_2 e^{\vk A} \|\PPPP_{k}^V{\mek}\|_{R_0}.
\end{equation}
To estimate $J_k$, we use the Markov property and~\eqref{5.4}
\begin{align*}
J_k(u)&\le A ^{-1} \E_u \left\{ \|u_k\|^2 \Xi_k^V\mm(u_k) \right\} \le A ^{-1}e^{\gamma}\, \E_u \left\{ \|u_k\|^2 \Xi_{k-1}^V \mm(u_k) \right\}\\&=A ^{-1} e^{\gamma}\, \E_u \left\{ \Xi_{k-1}^V\E_{u_{k-1}}\Big\{ \|u_1\|^2 \mm(u_1)\Big\} \right\} \le A^{-1} C_3 \PPPP_{k-1}^V\mm (u).
\end{align*}
Combining this with~\eqref{aaerfcvfd} and~\eqref{azzezfd}, and choosing $A$ so large that $q:=A^{-1} C_3<1$, we get
\[
\|\PPPP_k^V\mm\|_{L_\mm^\infty} \le C_2 e^{\vk A} \|\PPPP_{k}^V{\mek}\|_{R_0}+ q \|\PPPP_{k-1}^V\mm\|_{L_\mm^\infty}.
\]
Iterating, we obtain
\[
\|\PPPP_k^V\mm\|_{L_\mm^\infty} \le C_2 e^{\vk A} (1-q)^{-1} \|\PPPP_{k}^V{\mek}\|_{R_0}+ q^k.
\]
As $\PPPP_{k}^V{\mek}(u)\ge1$, we arrive at the required inequality~\eqref{dqdlkk}.
\end{step}

\begin{proof}[Step 2: Proof of~\eqref{estimE}]
For any $ t\ge1$, we have
\[
\PPPP_t^VF=\PPPP_{t-1}^V(\PPPP_1^VF)\le e^{\gamma}\,\PPPP_{t-1}^V(\PPPP_1F).
\]
So~\eqref{5.5} and~\eqref{dqdlkk} imply that
\[
\PPPP_t^VF(u) \le C_4\PPPP_{t-1}^V\wwww_8(u)\le C_5\PPPP_{t-1}^V\mm(u) \le C_6 \|\PPPP_{t}^V{\mek}\|_{R_0}\mm(u).
\]
This proves~\eqref{estimE}.
\end{proof}
\let\qed\relax
\end{proofc}

\subsection{Time-continuity}

The following lemma proves the time-continuity property.

\begin{lemm}
The function $t\mapsto\PPPP^V_tg(u) $ is continuous from~$\R_+$ to $\R$ for any $V\in C_b(H)$, $ g\in C_\wwww(H) $, $u\in H$, and $\wwww=\wwww_m$ with any $m\ge1$.
\end{lemm}

\begin{proof}
Let us show the continuity at the point $T\ge0$.
For any $t\ge0$, we write
\begin{align*}
\PPPP^V_Tg(u)-\PPPP^V_tg(u)&=\E_u \left\{\left[\Xi^V_T-\Xi^V_t \right] g(u_t)\right\}+ \E_u\left\{\left[g(u_T)- g(u_t) \right] \Xi^V_T\right\}\nonumber\\&=: S_1+S_2.
\end{align*}
As $V$ is bounded and $ g\in C_\wwww(H) $, we have
\begin{align*}
|S_1|&\le \E_u \left\{\left|\exp\left(\int_t^TV(u_s)\dd s\right)- 1 \right| \Xi^V_t|g(u_t)|\right\}\nonumber\\&\le \|g\|_{L^\ty_\wwww} \left(e^{|T-t|\|V\|_\infty}-1\right)e^{T\|V\|_\infty} \E_u \wwww(u_t).
\end{align*}
Combining this with~\eqref{5.3}, we get $S_1\to0$ as $t\to T$. To estimate $S_2$, we take any~$R>0$ and write
\begin{align*}
e^{-T\|V\|_\infty} | S_2|&\le \E_u \left| g(u_T)- g(u_t) \right| \\&= \E_u \left\{\I_{G_R^c} \left| g(u_T)- g(u_t) \right| \right\}+ \E_u \left\{\I_{G_R} \left| g(u_T)- g(u_t) \right| \right\}\\&=: S_3+S_4,
\end{align*}
where $G_R:=\{u_t, u_T\in B_H(R)\}$. From $ g\in C_\wwww(H) $ and~\eqref{5.3} we~derive
\begin{align*}
S_3 &\le C_1 \E_u \left\{\I_{G_R^c} (\wwww(u_T)+ \wwww(u_t)) \right\} \\&\le C_1 R^{-1} \E_u \left\{ \wwww^2(u_T)+ \wwww^2(u_t) \right\} \le C_2 R^{-1} \wwww^2(u).
\end{align*}
On the other hand, by the Lebesgue theorem on dominated convergence, for any $R>0$, we have $S_4 \to 0$ as $t\to T$. Choosing $R>0$ sufficiently large and $t$ sufficiently close to $T$, we see that $S_3+S_4$ can be made arbitrarily small. This shows that $S_2\to0$ as $t\to T$ and completes the proof of the lemma.
\end{proof}

\subsection{Uniform irreducibility}\label{S:2.3}

As $V$ is a bounded function, we have
\[
P_t^V(u,\dd v)\ge e^{-t\|V\|_\infty}P_t(u,\dd v),\quad\mbox{$u\in H$},
\]
where $P_t(u,{\cdot\,})$ is the transition function of the Markov family~$(u_t,\IP_u)$.
Thus to show the uniform irreducibility of $\{P^V_t\} $, it suffices to prove the following result.

\begin{prop}\label{P:2.7}
The family $\{P_t\} $ is uniformly irreducible with respect to the sequence $\{X_R\}$, i.e., for any integers $\rho, R\ge1$ and any $r>0$, there are positive numbers $l=l(\rho,r,R)$ and~$p=p(\rho,r)$ such that
\begin{equation}\label{3.4aa}
P_l(u,B_{H}(\hat u,r))\ge p,\quad\mbox{$u\in B_H(R),\,\hat u\in X_\rho$}.
\end{equation}
\end{prop}

\begin{proofc}

\begin{step}[Step~1]
There is a number $d>0$ such that for any $R\ge1$, we have
\begin{equation}\label{ezezrefs}
P_t(u, X_{d})\ge\frac12, \quad\mbox{$u\in B_H(R)$}
\end{equation}
for sufficiently large $t=t(R)$. Indeed, combining~\eqref{5.3}, \eqref{5.5}, and the Markov property, we get
\[
\E_u \|u_t\|_1^{2}\le C(e^{-8\alpha_1 t} R^{8}+1),\quad\mbox{$u\in B_H(R)$, $t\ge1$}.
\]
Taking $t$ so large that $e^{-8\alpha_1 t} R^{8}<1$ and $d >2\sqrt{ C} $ and using the Chebyshev inequality, we arrive at
\[
P_t(u, X_{d})\ge 1- d^{-2}C(e^{-8\alpha_1 t} R^{8}+1) \ge \frac12.
\]
\end{step}

\begin{proof}[Step~2]
By Lemma~3.3.11 in~\cite{KS-book}, for any non-degenerate ball $B \subset H$, there is~$p_1= p_1(d,B)>0$ such that
\[
P_1(u, B)\ge p_1,\quad\mbox{$u\in X_{d}$}.
\]
Combining this with a simple compactness and continuity argument, we get
\[
P_1(u, B_{H}(\hat u,r))\ge p_2,\quad\mbox{$u\in X_{d},\, \hat u\in X_\rho$}
\]
for some $p_2=p_2(d,\rho,r)>0$. This estimate, \eqref{ezezrefs}, and the Kolmogorov--Chapman relation imply~\eqref{3.4aa} with $l=t+1$ and $p=p_2/2$.
\end{proof}
\let\qed\relax
\end{proofc}

\subsection{Existence of an eigenvector}\label{S:2.4}

Here we show that the dual operator $\PPPP_t^{V*}$ has an eigenvector and give some decay estimates for it.

\begin{prop}\label{P:2.8}
For any $V\in C_b(H)$ and $t>0$, the operator~$\PPPP_t^{V*}$ has at least one eigenvector $\mu_{t,V}\in\PP(H)$ with a positive eigenvalue~$\lambda_{t,V}$:
\begin{equation}\label{6.22}
\PPPP_t^{V*}\mu_{t,V}=\lambda_{t,V}\mu_{t,V}.
\end{equation}
Moreover, any such eigenvector satisfies
\begin{align}
& \int_{H} \left(\|u\|_1^{n} + \mm_\vk (u) \right) \mu_{t,V}(\dd u)<\infty, \label{6.23}\\
&\|\PPPP_t^V\wwww_m\|_{ X_R }\int_{X_R^c}\wwww_m(u)\mu_{t,V}(\dd u)\to0\quad\text{as $R\to\infty$} \label{6.24}
\end{align}
for any $\vk\in (0,\gamma_0)$ and $n,m\ge1$.
\end{prop}

\begin{proofc}

\begin{step}[Step~1: Estimate~\eqref{6.23}]
Let us fix $t>0$, and let $\mu \in \PP(H)$ be an eigenvector of the operator $\PPPP_t^{V*}$ corresponding to an eigenvalue $\la>0$. Let us show that $\mu \in \PP_{\mm}(H)$ with $\mm=\mm_\vk$ for any $\vk\in (0,\gamma_0)$. Indeed, for any measurable function $f:H\to\R_+\cup\{+\ty\}$, we have
\begin{equation}
\lag f, \mu\rag=\lambda^{-1}\lag \PPPP_t^Vf, \mu\rag \le \lambda^{-1}e^{t\|V\|_\infty}\lag \PPPP_tf, \mu\rag.\label{6.32}
\end{equation}
Taking $f=\mm_\vk$, any number $A>0$, and setting\,\footnote{We do not indicate the dependence of different constants on $V,t,m,n,$ and $\vk$.} $C_1= \lambda^{-1}e^{t\|V\|_\infty}$, we obtain
\begin{align*}
\int_{H} \mm_\vk(u) \mu(\dd u)&\le C_1 \int_{H} \E_u \{\mm_{\vk}(u_t)\}\mu(\dd u) \\&=C_1 \!\!\int_{H}\!\!\Big(\E_u \left\{ \I_{\{\|u_t\|^2\le A\}} \mm_\varkappa(u_t) \right\}
\!+\! \E_u \left\{ \I_{\{\|u_t\|^2> A\}} \mm_\varkappa(u_t) \right\}\!\Big)\mu(\dd u)\\&\le C_1 \int_{H}\big(\exp(\vk A)+A^{-1} \E_u \left\{ \|u_t\|^2 \mm_\varkappa(u_t) \right\}\big)\mu(\dd u)\\&\le C_1 \int_{H}\big(\exp(\vk A)+C_2A^{-1} \mm_\vk(u)\big)\mu(\dd u),
\end{align*}
where we used inequality~\eqref{5.4}. Choosing $A>C_1C_2$, we get
\begin{equation}\label{sqddsf}
\int_{H} \mm_\vk(u) \mu(\dd u)\le C_1(1- C_1C_2 A^{-1})^{-1}\exp(\vk A)<\ty,
\end{equation}
so\,\footnote{Note that this proof is formal. A rigorous proof can be obtained by applying the above arguments to bounded approximations of~$\mm$.} $\mu\in \PP_\mm(H)$. Taking $f(u)=\|u\|_1^n $ in~\eqref{6.32} and using~\eqref{5.5} and~\eqref{sqddsf}, we~obtain
\[
\int_{H} \|u\|_1^n \mu(\dd u) \le C_1 \int_{H} \E_u \{\|u_t\|_1^n\}\mu(\dd u) \le C_3 \int_{H} (1+\|u\|^{8n})\mu(\dd u)<\ty
\]
for any $n\ge1$. This proves~\eqref{6.23}.
\end{step}

\begin{step}[Step~2: Limit~\eqref{6.24}]
From~\eqref{5.3} it follows that
\begin{align}\label{6.036}
\|\PPPP_t^V\wwww_m\|_{X_R}&\le e^{t\|V\|_\infty}\sup_{u\in X_R}\E_u\wwww_m(u_t)\nonumber\\&
\le C_4 \,\sup_{u\in B_H(R)}\wwww_m(u)= C_4(1+ R^{2m}).
\end{align}
Using the Cauchy--Schwarz inequality, \eqref{6.23}, and the Chebyshev inequality, we see that
\[
\int_{X_R^c}\wwww_m(u)\mu(\dd u)
\le \langle\wwww_m^2,\mu\rangle^{1/2} \mu\bigl(X_R^c\bigr)^{1/2}\le C_{5}R^{-n}.
\]
Combining this with~\eqref{6.036} and choosing $n>2m$, we obtain~\eqref{6.24}.
\end{step}

\begin{proof}[Step~3: Construction of an eigenvector]
Let us take any $A>0$ and~$m\ge 1$ and define the convex set
\[
D_{A,m}:=\bigl\{\nu\in\PP(H): \langle\wwww_m,\nu\rangle\le A\bigr\}.
\]
By the Fatou lemma, $D_{A,m}$ is closed in~$\PP(H)$.
Consider the continuous mapping
\[
G:=G(t,V):D_{A,m}\to \PP(H), \quad\quad \nu \mapsto \frac{\PPPP_t^{V*}\nu}{\PPPP_t^{V*}\nu(H)}.
\]
Let us show that $G(D_{A,m})\subset D_{A,m}$ for an appropriate choice of~$A$ and~$m$, and that~$G(D_{A,m})$ is compact in~$\PP(H)$. In view of the Leray--Schauder theorem, this will imply the existence of an eigenvector~$\mu\in D_{A,m}$ satisfying~\eqref{6.22} with eigenvalue $\la=\PPPP_t^{V*}\mu(H)>0$. From~\eqref{5.3} we derive~that
\begin{align*}
\bigl\langle\wwww_m,G(\nu)\bigr\rangle &\le \exp\{t\Osc(V)\}\lag \wwww_m, \PPPP_t^*\nu \rag
\notag\\
&\le \exp\{t(\Osc(V)-m\alpha_1)\}\langle\wwww_m,\nu\rangle+C_{6},
\end{align*}
where $ \Osc(V):=\sup_{u\in H} V(u)-\inf_{u\in H} V(u)$ is the oscillation of $V$. Choosing~$A$ and~$m$ so large that $\exp\{t(\Osc(V)-m\alpha_1)\}\le 1/2$ and $A\ge 2 C_{6}$, we get that $G(D_{A,m})\subset D_{A,m}$. In view of the Prokhorov compactness criterion (see~\cite[Theorem~11.5.4]{dudley2002}), to prove that~$G(D_{A,m})$ is relatively compact, it suffices to check that
\[
\int_{H}\|u\|_1^2 \PPPP_t^{V*}\nu(\dd u)\le C_{7}\quad\mbox{for any $\nu\in D_{A,m}$}.
\]
Using~\eqref{5.5} and the fact that $V$ is bounded, we get
\begin{align*}
\int_{H}\|u\|_1^2 \PPPP_t^{V*}\nu(\dd u)&\le \exp(t\|V\|_\ty)\int_{H}\|u\|_1^2 (\PPPP_t^*\nu)(\dd u)\\&\le C_8 \int_H\|u\|^8 \nu(\dd u) \\&\le C_{9} \int_H\wwww_m(u) \nu(\dd u) \le C_{9} A=:C_7.
\end{align*}
Thus there is an eigenvector $\mu\in D_{A,m}$.
\end{proof}
\let\qed\relax
\end{proofc}


\section{Uniform Feller property}\label{S:3}

In this section, we establish the following result.

\begin{theo}\label{T:3.1}
For any $V\in \VV$, the family $\{P_t^V\} $ satisfies the uniform Feller property with respect to the sequence $\{X_R\}$, i.e., there is an integer $R_0\ge1$ such that the family $\{\|\PPPP_t^V{\mek}\|_R^{-1}\PPPP_t^V\psi,t\ge0\}$ is uniformly equicontinuous on~$X_R$ for any~$\psi\in \VV$ and~$R\ge R_0$.
\end{theo}

See the papers~\cite{JNPS-2012, JNPS-2014, MN-2017} for similar results in the case of a discrete-time random dynamical system and~\cite{MN-2015} for the case of the stochastic damped nonlinear wave~equation. The main difficulty in the proof of Theorem~\ref{T:3.1} comes from the fact that the oscillation of the potential $V$ can be arbitrarily large. To overcome this, we introduce a new auxiliary equation in the construction of the coupling processes and choose carefully the parameters in order to have a stabilisation property with an appropriate rate.

\subsection{Construction of coupling processes} \label{S:3.1}

The coupling processes are constructed following the arguments of~\cite{MN-2015}. Let us take any~$z,z' \in H$ and denote by $u_t$~and~$u_t'$ the solutions of~\eqref{0.4} issued from $z$ and~$z'$. For any integer $N\ge 1$ and number $\la>0$, let $v$ be the solution of the following problem
\begin{equation}\label{auxi}
\dot v + B(v) + Lv + {\mathsf P}_N[\la(v-u)+B(u)-B(v)] = h + \eta(t), \quad v(0)=z',
\end{equation}
where $\eta$ is defined by~\eqref{0.2}. We denote by~$\nu(z,z')$ and~$\nu'(z')$ the laws of processes~$\{v(t), t\in J\}$ and $\{u'(t), t\in J\}$, respectively, where $J=[0,1]$. We shall use the following result.

\begin{prop}\label{P:dist}
There exists an integer $N_1\ge1$ such that if $N\ge N_1$ and~$\la\ge N^2/2 $, then for any~$\e>0$ and $z,z'\in H$, we have
\begin{equation}\label{eoejtnvf}
\| \nu(z,z')-\nu'(z')\|_{var}\le \e^a+ 2 \left[\exp\left(C_{\la,N} \e^{a-2}\|z-z'\|^2 e^{C(\|z\|^2+\|z'\|^2)} \right)-1\right]^{1/2},
\end{equation}
where $\|{\,\cdot\,}\|_{var}$ denotes the total variation distance on $\PP(C(J; H))$ and $a<2, C,$ and~$C_{\la,N}$ are positive constants not depending on $\e, z,z'$.
\end{prop}

See Section~\ref{apascss} for the proof. By Proposition~1.2.28 in~\cite{KS-book}, there is a~probability space $(\hat \Omega, \hat \FF, \hat\pP)$ and measurable functions $\mathcal{Z}, \mathcal{Z}' :H\times H\times \hat \Omega\to C(J;H) $ such that~$(\mathcal{Z} (z,z'), \mathcal{Z}'(z,z'))$ is a maximal coupling for~$(\nu(z,z'), \nu(z'))$ for any $z,z'\in H$. We denote by~$\tilde v$ and $\tilde u'_t$ the restrictions of $\mathcal{Z}$ and $ \mathcal{Z}' $ to time $t\in J$. Then $\tilde v_t$ is a solution of
\[
\dot {\tilde v} + B(\tilde v) + L\tilde v + {\mathsf P}_N[\la \tilde v-B(\tilde v)] = h +\psi(t), \quad \tilde v (0)= z',
\]
where the process $\{\int_0^t\psi(s)\dd s,\, t\in J\}$ has the same law as
\[
\left\{W(t)- \int_0^t {\mathsf P}_N[B(u_s)-\la u_s]\dd s, \,\, t\in J\right\}.
\]
Let $\tilde u_t$ be a solution of
\[
\dot { \tilde u}+B(\tilde u)+L \tilde u + {\mathsf P}_N[\la \tilde u-B(\tilde u)]=h +\psi(t), \quad \tilde u (0)= z.
\]
Then $\{\tilde u_t, t\in J\}$ has the same law as $\{ u_t, t\in J\}$. Now the coupling operators~$\RR$ and~$\RR'$ are defined by
\[
\RR_t(z,z',\omega)=\tilde u_t,\quad \RR'_t(z,z',\omega)=\tilde u_t', \quad z, z'\in H,\, \omega\in\hat\Omega,\, t\in J.
\]
By Proposition~\ref{P:dist}, for any $\e>0$, $N\ge N_1$, and $\la\ge N^2/2 $, we have
\begin{multline}\label{sdsflklk}
\hat \pP\{\exists t\in J\text{ s.t. } \tilde v_t\neq \tilde u_t' \}\\
\le \e^a+ 2 \left[\exp\left(C_{\la,N} \e^{a-2}\|z-z'\|^2 e^{C(\|z\|^2+\|z'\|^2)} \right)-1\right]^{1/2}.
\end{multline}
Let $(\Omega^k,\FF^k,\pP^k), k\ge0$ be a sequence of independent copies of $(\hat \Omega, \hat \FF, \hat\pP)$ and $(\Omega,\FF,\pP)$ the direct product of $(\Omega^k,\FF^k,\pP^k)$. For any $\omega=(\omega^1,\omega^2,\ldots)\in\Omega$ and $z,z'\in H$, we set~$\tilde u_0=z$, $\tilde u_0'=z'$, and
\begin{align*}
\tilde u_t(\omega)&=\RR_s(\tilde u_{k}(\omega),\tilde u'_{k}(\omega),\omega^k), &\quad
\tilde u'_t(\omega)&=\RR_s'(\tilde u_{k}(\omega),\tilde u'_{k}(\omega),\omega^k), \\\tilde v_t(\om)&=\mathcal{Z}_s(\tilde u_{k}(\omega),\tilde u'_{k}(\omega),\omega^k),
\end{align*}
where $t=s+k, s\in [0,1)$. We shall say that $(\tilde u_t,\tilde u_t')$ is a \emph{coupled trajectory at level}~$(N,\la)$ issued from~$(z,z')$.

\subsection{Proof of Theorem~\ref{T:3.1}}

\begin{step}[Step~1:~Stratification]
Let us take any functions $V,\psi\in \VV$ and points $ z,z' \in X_R$ such that $d:=\|z-z'\|\le1$. We need to prove the uniform equicontinuity of the family $\{g_t,t\ge 0\}$ on~$X_R$, where
\[
g_t= \|\PPPP_t^V{\mek}\|_R^{-1}\PPPP_t^V\psi.
\]
Without loss of generality, we can assume that $\psi$ and $V$ are non-negative,~$ \psi\le~1$, and the integer $N$ in representation~\eqref{repres} is the same for $\psi$ and $V$ (we denote it by~$N_0$). Let $(u_t,u_t'):=(\tilde u_t,\tilde u_t')$ be a coupled trajectory~at level~$(N,\la)$ issued from~$(z,z')$ and let~$v_t:=\tilde v_t$ be the associated process. The parameters $N\ge N_0$ and~$\lambda\ge N^2/2$ will be chosen later.

%\smallskip
Following~\cite{MN-2015,JNPS-2014}, for any integers $r\ge0$ and~$\rho\ge1$, we introduce the~events\,\footnote{The event $\bar G_r$ is well defined also for $r=+\ty$.}
\begin{align*}
\bar G_r &=\bigcap_{j=0}^rG_j, \,\,\, G_j=\{ v_t =u_t', \forall t\in (j,j+1]\}, \,\,\, F_{r,0}=\varnothing, \\
F_{r,\rho}&=\biggl\{ \sup_{t \in [0,r]} \left(\int_0^t \left(\| u_s \|_1^2+\| u'_s\|_1^2\right)\dd s-Kt \right) \le \|z\|^2+ \|z'\|^2 + \rho; \\
& \mkern 400mu \|u_r\|^2+ \|u_r'\|^2 \le \rho \biggr\},
\end{align*}
where $K$ is the constant in~\eqref{5.6}, and the pairwise disjoint events
\[
A_0=G_0^c, \,\,\, A_{r,\rho}=\bigl(\bar G_{r-1}\cap G_r^c\cap F_{r,\rho}\bigr)\setminus F_{r,\rho-1},\,\, r\ge1,\,\rho\ge1, \,\,\, \tilde A = \bar G_{+\ty}
\]
We decompose as follows
\begin{align}
\PPPP_t^V\psi(z)-\PPPP_t^V\psi(z') &= \E \bigl\{\I_{A_{0}}\bigl[\Xi^V_t\psi (u_t) -\Xi^V_t\psi (u_t')\bigr] \bigr\}\nonumber\\
&\qquad+ \sum_{r,\rho=1}^{\ty} \E \bigl\{\I_{A_{r,\rho}}\bigl[\Xi^V_t\psi (u_t) -\Xi^V_t\psi (u_t')\bigr] \bigr\}\nonumber\\
&\qquad+ \E \bigl\{\I_{\tilde A}\bigl[\Xi^V_t\psi (u_t) -\Xi^V_t\psi (u_t')\bigr] \bigr\}\nonumber\\
&=I_0^t+ \sum_{r,\rho=1}^{\ty} I_{r,\rho}^{t}+ \tilde I^{t},\label{E:7.3}
\end{align}
where
\begin{align*}
I_{0}^{t}&=\E \bigl\{\I_{A_{0}}[\Xi^V_t\psi (u_t) -\Xi^V_t\psi (u_t')\bigr] \bigr\},\\
I_{r,\rho}^{t}&=\E \bigl\{\I_{A_{r,\rho}}[\Xi^V_t\psi (u_t) -\Xi^V_t\psi (u_t')\bigr] \bigr\},\\
\tilde I^{t}&=\E \bigl\{\I_{\tilde A}\bigl[\Xi^V_t\psi (u_t) -\Xi^V_t\psi (u_t')\bigr] \bigr\}.
\end{align*}
\end{step}

In Steps 2 and 3, we estimate $I_{0}^t,$ $I_{r,\rho}^t$, and $\tilde I^{t}$.

\begin{step}[Step~2:~Estimates for $I_{0}^t$ and $I_{r,\rho}^t$]
We have following inequalities
\begin{align}
|I_{0}^t|&\le C_1(R,V)\|\PPPP_t^V{\mathbf1}\|_R\,\pP(A_{0})^{1/2}, \label{6.5aaa}
\\
|I_{r,\rho}^t|&\le C_2(R,V)e^{r\|V\|_\infty}\|\PPPP_t^V{\mathbf1}\|_R\,\pP(A_{r,\rho})^{1/2} \label{6.5}
\end{align}
for any integers $r,\rho\ge1$ and $R\ge R_0$, where $R_0$ is the number in Proposition~\ref{P:2.4}. Let us prove~\eqref{6.5}, the other inequality is proved in a similar way. First assume that~$r+1\le t$. Using $ \psi \le 1$, the positivity of~$\Xi^V_t\psi $, and the Markov property, we~derive
\begin{align*}
I_{r,\rho}^t &\le\E\bigl\{\I_{A_{r,\rho}}\Xi^V_t\psi(u_t)\bigr\}
\le \E\bigl\{\I_{A_{r,\rho}}\Xi^V_t\bigr\}\\
&= \E\bigl\{\I_{A_{r,\rho}}\E\bigl[\Xi^V_t\,\big|\,\FF_{r+1}\bigr]\bigr\}
\le e^{r\|V\|_\infty}\E\bigl\{\I_{A_{r,\rho}}(\PPPP_{t-r-1}^V{\mathbf1})(u_{r+1})\bigr\},
\end{align*}
where $\{\FF_t\}$ is the filtration generated by~$(u_t,u_t')$. Then from the positivity of~$V$ and~\eqref{a5.8} it follows~that
\[
\PPPP_{t-r-1}^{V}{\mathbf1}(y)\le \PPPP_{t}^{V}{\mathbf1}(y)\le M\|\PPPP_{t}^{V}{\mathbf1}\|_{R_0}\wwww (y),\quad y\in H,
\]
so that
\begin{align*}
I_{r,\rho}^t &\le C_3 e^{r\|V\|_\infty}\|\PPPP_{t}^V{\mathbf1}\|_{R_0}\E\bigl\{\I_{A_{r,\rho}}\wwww (u_r)\bigr\}\\
&\le C_3 e^{r\|V\|_\infty}\|\PPPP_{t}^V{\mathbf1}\|_{R_0}\bigl\{\pP(A_{r,\rho})\,\E\,\wwww ^2(u_r)\bigr\}^{1/2}.
\end{align*}
Using this,~\eqref{5.3}, and the symmetry, we obtain~\eqref{6.5}. If $r> t$, then
\begin{align*}
I_{r,\rho}^t &
\le e^{r\|V\|_\infty}\pP(A_{r,\rho})\le e^{r\|V\|_\infty} \|\PPPP_t^V{\mathbf1}\|_R\,\pP(A_{r,\rho})^{1/2},
\end{align*}
which implies~\eqref{6.5} by symmetry.
\end{step}

\begin{step}[Step~3:~Estimate for $\tilde I^t$]
Let us show that
\begin{equation}\label{E:eerrtt}
|\tilde I^{t}|\le C_{4}(R,V,\la,N, \psi) \|\PPPP_t^V{\mathbf1}\|_R d.
\end{equation}
Indeed, we write
\[
\tilde I^t =\E\bigl\{\I_{\tilde A}\Xi^V_t[\psi (u_t)-\psi (u_t')] \bigr\}+\E\bigl\{\I_{\tilde A}[\Xi^V_t-\Xi^{V'}_t]\psi (u_t') \bigr\} =:J_1^t+J_2^t,
\]
where $\Xi^{V'}_t:=\exp\left(\int_{0}^tV(u_s') \dd s\right)$. Then by~\eqref{FPE1}, on the event $\tilde A$ we have
\[
\|\pp(u_s-u_s')\|
\le e^{-\lambda s}d, \quad s\in [0,t].
\]
Since $\psi\in L_b(H) $, we derive from this
\begin{align*}
|J_{1}^t|&\le \E\bigl\{\I_{\tilde A} \Xi^V_t| \psi (u_t)-\psi (u_t')| \bigr\} \le \| \psi\|_{L} e^{-\lambda t} \|\PPPP_t^V{\mathbf1}\|_R d \le \| \psi\|_{L} \|\PPPP_t^V{\mathbf1}\|_Rd.
\end{align*}
Similarly, as $V\in L_b(H)$,
\begin{align*}
|J_{2}^t| \le \E\bigl\{\I_{\tilde A}|\Xi^V_t-\Xi^{V'}_t| \bigr\} &\le \E\left\{\I_{\tilde A}\Xi^V_t \left[\exp\left(\int_{0}^t|V(u_s)-V(u_s')| \dd s\right)-1 \right] \right\}\\&\le \left[\exp\left(C_{\la,N} \lambda^{-1} \|V\|_{L} d (1-e^{-\lambda t})\right)-1 \right] \|\PPPP_t^V{\mathbf1}\|_R\\&\le \left[\exp\left(C_5(R,V,\la,N) d \right)-1 \right] \|\PPPP_t^V{\mathbf1}\|_R.
\end{align*}
Recalling that $d\le1$ and combining the estimates for $J_{1}^t$ and $J_{2}^t$, we get~\eqref{E:eerrtt}.
\end{step}

\begin{step}[Step~4:~Uniform equicontinuity of $g_t$]
We use the following lemma, which is proved at the end of this subsection.

\begin{lemm}\label{L:4.3PR}
For any $\alpha>0$,
there is an integer $N_2(\alpha)\ge1$ and positive numbers $a$ and $\beta$ such~that
\begin{align}
\pP\{A_{0}\} &\le C_6(R,\la, N)d^{a/2},\label{E:7.7ffd}\\
\pP\{A_{r,\rho}\} &\le C_7(R) \!\left\{ \left(d^ae^{-a\alpha r} \!+\! \left[\exp \left(C_8(R,\la,N)d^a e^{ C'\rho-a\alpha r} \right) \!-\!1\right]^{1/2} \!\right) \!\wedge \! e^{-\beta\rho} \right\} \label{E:7.7}
\end{align}
for any $N\ge N_2(\alpha)$, $\la\ge N^2/2 $, $R\ge1$, and a universal constant $C'>0$.
\end{lemm}

From~\eqref{E:7.3}--\eqref{E:7.7} it follows that, for any $z,z'\in X_R$, $t\ge0$, $R\ge R_0$, and~$\alpha>0$, we have
\begin{multline*}
\bigl|g_t(z)-g_t(z')\bigr|\le C_{9}(R,V,\la,N, \psi) \Bigg(d^{a/4} + d \\
+ \sum_{r,\rho=1}^\infty \!\! e^{r\|V\|_\infty} \! \left\{ \! \left(d^{a/2}e^{-a\alpha r/2} + \left[\exp \left(C_8d^a e^{C'\rho-a\alpha r} \! \right) -1\right]^{1/4} \right) \wedge e^{-\beta\rho/2} \right\} \!\Bigg),
\end{multline*}
provided that $N \ge N_0\vee N_1\vee N_2(\alpha) $ and $\la\ge N^2/2 $. When $d=0$, the series on the right-hand side vanishes. So to prove the uniform equicontinuity of~$\{g_t\}$, it suffices to show that the series converges uniformly in~$d\in[0,1]$. Since its terms are positive and monotone, it suffices to show the converge for~$d=1$:
\begin{align}
\sum_{r,\rho=1}^\infty e^{r\|V\|_\infty} \left\{ \left(e^{-a\alpha r/2} + \left[\exp \left(C_8 e^{C'\rho-a\alpha r} \right) -1\right]^{1/4} \right) \wedge e^{-\beta\rho/2} \right\} < \infty. \label{E:7.5}
\end{align}
To prove this, we will assume that~$\alpha$ is sufficiently large. Let
\[
S_1=\{(r,\rho)\in\N^2: \rho\le a\alpha r/(2C') \}, \quad S_2=\N^2\setminus S_1.
\]
Then taking $\alpha> 16 \|V\|_\ty/a$, we see that
\begin{multline*}
\sum_{(r,\rho)\in S_1}e^{r\|V\|_\infty} \left(e^{-a\alpha r/2} + \left[\exp \left(C_8 e^{C'\rho-a\alpha r} \right) -1\right]^{1/4} \right) \\
\le C_{10}(R,N) \sum_{(r,\rho)\in S_1} e^{r\|V\|_\infty} e^{-a\alpha r/8}\le C_{11}(R,N) \sum_{r=1}^\ty e^{-a\alpha r/16}<\infty.
\end{multline*}
Choosing $\alpha> 8C'\|V\|_\ty/(a\beta) $, we get
\[
\sum_{(r,\rho)\in S_2} e^{r\|V\|_\ty}e^{-\beta\rho/2}\le C_{12} \sum_{\rho=1}^\infty e^{-\beta\rho/4}<\infty.
\]
These two inequalities show that~\eqref{E:7.5} holds.

\begin{proof}[Proof of Lemma~\ref{L:4.3PR}]
Taking $\e=d$ in~\eqref{sdsflklk} and using $d\le1$, we get
\[
\IP\bigl\{ A_0 \bigr\}
\le d^a+2\left[\exp\left(C_{\la, N} d^{a} e^{ {2C R^2}} \right)-1\right]^{1/2}\le C_6(R,\la, N)d^{a/2}
\]
for $N \ge N_1$ and $\la\ge N^2/2 $.
This gives~\eqref{E:7.7ffd}. From the inclusion $A_{r,\rho}\subset F_{r,\rho-1}^c$ and inequalities~\eqref{5.6} and~\eqref{5.2} it follows that
\begin{equation}\label{6.7}
\pP\{A_{r,\rho}\}\le C_{13}(R) e^{-\beta\rho},
\end{equation}
where $ \beta:= \gamma_0/2$. By Proposition~\ref{4.13}, on the event $A_{r,\rho}$ we have
\begin{equation}
\|u_{r}-u_{r}'\|
\le \exp(-\alpha r + c(\| z\|^2 +\|z'\|^2+\rho))d\le C_{13}(R) e^{-\alpha r +c\rho}d, \label{FP:est}
\end{equation}
provided that $N\ge N_1'(\alpha):=\sqrt{\alpha +cK} $ and $\la\ge N^2/2 $. Recall that on the same event we have also
\begin{equation}\label{FP:est1}
\| u_r\|^2 +\|u_r'\|^2\le \rho.
\end{equation}
Using the Markov property, \eqref{sdsflklk} with $\e=de^{-\alpha r}$, \eqref{FP:est} and~\eqref{FP:est1},~we~obtain
\begin{align*}
&\IP\{A_{r,\rho}\}\le\IP\bigl\{\bar G_{r-1} \cap G_r^c \cap F_{r,\rho}\bigr\} =\E\bigl\{\I_{\bar G_{r-1}\cap F_{r,\rho}}\E\bigl(\I_{G_r^c}\,\big|\,\FF_{r}\bigr)\bigr\}\\
&\le d^ae^{-a\alpha r} +2\E\Big\{\I_{\bar G_{r-1}\cap F_{r,\rho}}\!\! \left[\exp\left(C_{\la,N} d^{a-2} e^{-(a-2)\alpha r}\|u_r\!-\!u_r'\|^2 e^{C(\| u_r\|+\|u_r'\|)} \right)-1\right]^{1/2}\!\Big\}\\
&\le d^ae^{-a\alpha r}+2\left[\exp\left(C_8(R,\la,N)d^a e^{C'\rho-a\alpha r} \right)-1\right]^{1/2}.
\end{align*}
Combining this with~\eqref{6.7} and taking $N\ge N_2(\alpha):=N_1\vee N_1'(\alpha) $ and $\la\ge N^2/2 $, we get the required inequality~\eqref{E:7.7}.
\end{proof}
\end{step}


\section{Proof of Theorem~\ref{T:1.1}}\label{S:3Sd}

The results of Sections~\ref{S:2} and~\ref{S:3} show that the conditions of Theorem~\ref{T:5.3} are satisfied if we choose
\begin{gather*}
P_t^V(u,\Gamma)=(\PPPP_t^{V*} \delta_u) (\Gamma),\,\,\,X=H, \,\,\, X_R= B_{H^1}(R), \,\, R\ge R_0, \,\,\, \\
\wwww(u)=\wwww_m(u)=1+\|u\|^{2m},\,\,\, \CC= \VV, \,\,\, V\in \VV
\end{gather*}
with sufficiently large $m$ and $R_0$.~Thus there are eigenvectors $\mu_V\in \PP(H)$ and $h_V\in L_\wwww^\infty(H)$ corresponding to an eigenvalue~$\lambda_V>0$.~Moreover, for~any~$R\ge1$, the~restriction~of~$h_V$ to~$X_R$ is continuous and strictly positive, so~$h_V: H^1\to \R$ is continuous and strictly positive.~As $\pP_u\{u_1\in H^1\}=1$ and $h(u)=\la_V^{-1}\PPPP_1^Vh_V(u)$, we have
\[
h_V(u)\ge\lambda_V^{-1}e^{- \|V\|_\infty}\,\E_u h_V(u_{1})>0 \quad u\in H.
\]
The continuity of $h_V:H\to \R $ follows from the uniform convergence in~\eqref{1.5}, and the uniqueness of $\mu_V$ and~$h_V$ from~\eqref{1.5} and~\eqref{1.6}. The proof of~\eqref{1.5} is carried out in Steps~1--3, and that of~\eqref{1.7} in Step 4. Convergence~\eqref{1.6} follows immediately from~\eqref{1.5}.

\begin{step}[Step~1: Proof of~\eqref{1.5} for $f\in \VV$]
In view of~\eqref{5.12}, for any $f\in \VV$, we have limit~\eqref{1.5} in $C(X_R)\cap L^1(H,\mu_V)$. We claim that this limit holds also in~$C(B_H(R))$ for any $R\ge1$.~Indeed, it suffices to check condition~\eqref{5.13} with~$B=B_H(R)$ and~$s=1$, i.e.,
\[
A_{R,r}:= \sup_{u\in B_H(R)} \int_{H\setminus X_r}\wwww_m(v)\,P^V_1(u,\dd v) \to0
\quad\mbox{as $r\to\infty$}.
\]
From the Poincar\'e inequality and~\eqref{5.5} it follows that
\begin{align*}
A_{R,r}&\le r^{-2}\sup_{u\in B_H(R)} \E_u\left\{\wwww_m(u_1)\|u_1\|_1^2 \,\Xi_1^V \right\}\\&\le r^{-2}e^{\|V\|_\ty}\sup_{u\in B_H(R)} \E_u\left\{(1+\alpha_1^{-m}\|u_1\|_1^{2m})\|u_1\|_1^2 \right\} \\&\le r^{-2} C(m) R^{8(m+1)}\to 0\quad\text{as $r\to\ty$}.
\end{align*}
This implies~\eqref{1.5} for $f\in \VV$.
\end{step}

\begin{step}[Step~2: Proof of~\eqref{1.5} for $f\in C_b(H)$]
For any $n\ge1$, let $\tilde f_n\in L_b(H)$ be such that
\[
\sup_{u\in B_H(n)}|\tilde f_n(u)-f(u)|\le\frac1n.
\]
Then the functions $f_n={\tilde f}_n\circ{\mathsf P}_n$ belong to the space~$\VV$, satisfy $\|f_n\|_\infty\le \|f\|_\infty$ and~$f_n\to f$ as $n\to\infty$, uniformly on compact subsets of~$H$. Setting
\[
\Delta_t(g)=\sup_{u\in B_H(R)}\bigl|\lambda_V^{-t}\PPPP_t^Vg(u)-\langle g,\mu_V\rangle h_V(u)\bigr|,
\quad \|g\|_{0,R}=\sup_{u\in B_H(R)}|g(u)|,
\]
for any $t\ge0$ and $n\ge1$, we write
\[
\Delta_t(f)\le \Delta_t(f_n)+\|h_V\|_{0,R}\,|\langle f-f_n,\mu_V\rangle|+\lambda_V^{-t}\|\PPPP_t^V(f-f_n)\|_{0,R}.
\]
Since $f_n\in\VV$, the first term on the right-hand side of this inequality goes to zero as $k\to\infty$ for any fixed $n\ge1$. The Lebesgue theorem on dominated convergence implies that $|\langle f-f_n,\mu_V\rangle|\to0$ as $n\to\infty$. Thus, the required~convergence will be established if we show that
\begin{equation}\label{6.57EE}
\sup_{t\ge 1}\lambda_V^{-t}\|\PPPP_t^V(f-f_n)\|_{0,R}\to0\quad\mbox{as $n\to\infty$}.
\end{equation}
To prove~this limit, we take any $\rho>0$ and write
\begin{equation}\label{6.58}
\|\PPPP_t^V(f-f_n)\|_{0,R}\le J_1(t,n,\rho)+J_2(t,n,\rho),
\end{equation}
where
\[
J_1(t,n,\rho)=\|\PPPP_t^V\bigl((f-f_n)\I_{X_\rho}\bigr)\bigr\|_{0,R}, \quad J_2(t,n,\rho)=\|\PPPP_t^V\bigl((f-f_n)\I_{X_\rho^c}\bigr)\|_{0,R}.
\]
By~\eqref{6.0015}, we have
\[
J_1(t,n,\rho)\le \e(n,\rho)\,\|\PPPP_t^V\mathbf1\|_{0,R} \le \e(n,\rho)\,C_R \,\|\PPPP_t^V\mathbf1\|_{R_0},
\]
where $\e(n,\rho)=\|f-f_n\|_{X_\rho}\to0$ as $n\to\infty$. Convergence~\eqref{1.5} with $f=\mathbf1$ implies that
\begin{equation}\label{VVN}
\text{the set $\{\lambda_V^{-t}\|\PPPP_t^V{\mek }\|_{R_0}\}_{t\ge0}$ is bounded in $\R$.}
\end{equation}
It follows that
\begin{equation}\label{6.59}
\sup_{t\ge 0}\lambda_V^{-t}J_1(t,n,\rho) \to0\quad\mbox{as $n\to\infty$}.
\end{equation}
To estimate $J_2$, we use~\eqref{estimE}. For any $ \rho,n \ge1$ and $t\ge0$, we have
\begin{align*}
\la_V^{-t}J_2(t,n,\rho) &\le 2\|f\|_\infty \rho ^{-2} \la_V^{-t} \|\PPPP_t^V F \|_{0,R}\\
&\le C_R\|f\|_\infty \rho ^{-2} \la_V^{-t}\|\PPPP_t^V{\mathbf1}\|_{R_0}.
\end{align*}
By~\eqref{VVN}, the right-hand side of this inequality goes to zero as $\rho\to\infty$, uniformly with respect to $t\ge1$. Combining this with~\eqref{6.59}, we see that supremum over $t\ge 1$ of the right-hand side of~\eqref{6.58} can be made arbitrarily small by choosing first $\rho>0$ and then~$n\ge1$ sufficiently large. This proves~\eqref{6.57EE}.
\end{step}

%\smallskip
\begin{step}[Step~3: Proof of~\eqref{1.5} for $f\in C_\mm(H)$]
We use again an approximation~argument. Let us fix any $\vk\in(0,\gamma_0)$ and $f\in C_{\mm}(H)$ with $\mm=\mm_\vk$. We define a sequence~$\{f_n\}$ by the relation $f_n=f^+\wedge n-f^-\wedge n$. Then $f_n\in C_b(H)$, $|f_n|\le |f|$ for any $n\ge1$, and $f_n\to f$ in $L_{\mm'}^\infty(H)$ with $\mm'=\mm_{\vk'}$ for any $\vk'\in(\vk,\gamma_0)$. Furthermore, in view of~\eqref{1.5} and the Lebesgue theorem on dominated convergence, we~have
\begin{gather*}
\sup_{u\in B_H(R)}\bigl|\lambda_V^{-t}\PPPP_t^Vf_n(u)-\langle f_n,\mu_V\rangle h_V(u)\bigr| \to0\quad\mbox{as $t\to\infty$ for any fixed $n\ge1$},\\
|\langle f-f_n,\mu_V\rangle|\to0\quad\mbox{as $n\to\infty$}.
\end{gather*}
Thus, as in the previous step, it suffices to prove that
\begin{equation}\label{6.57}
\sup_{t\ge 0}\lambda_V^{-t}\|\PPPP_t^V(f-f_n)\|_{0,R}\to0\quad\mbox{as $n\to\infty$}
\end{equation}
To see this, we use~\eqref{dqdlkk} for $\mm'$:
\[
\|\PPPP_t^V(f-f_n)\|_{0,R}\le \e_n\,\|\PPPP_t^V\mm'\|_{0,R}
\le C_R\,\e_n\, \|\PPPP_t^V{\mek}\|_{R_0},
\]
where $\e_n=\|f-f_n\|_{L_{\wwww'}^\infty}\to0$ as $n\to\infty$. Combining this with~\eqref{VVN}, we get~\eqref{6.57}.
\end{step}

\begin{step}[Step~4: Proof of~\eqref{1.7}]
In view of~\eqref{1.5}, it suffices to show that
\begin{equation}\label{6.63}
\sup_{(t, \nu)\in\R_+\times\Lambda({\vk'},M)}\biggl\{\int_{B_H(R)^c}\bigl|\lambda_V^{-t}\PPPP_t^Vf-\langle f,\mu_V\rangle h_V\bigr|\,\nu(\dd u) \biggr\}\to0\,\,\,
\mbox{as $R\to\infty$}.
\end{equation}
By~\eqref{dqdlkk} and~\eqref{VVN}, we have
\[
\|\PPPP_t^Vf\|_{L_\mm^\infty}\le C_1\|\PPPP_t^V{\mek}\|_{R_0}\le C_2\lambda_V^t\quad\mbox{for all $t\ge0$}.
\]
It follows that
\[
\bigl|\lambda_V^{-k}\PPPP_t^Vf(u)\bigr|\le C_3\mm_\vk(u).
\]
Since $ \vk< \vk'$, $h_V\in C_\wwww(H)$, and
\[
\sup_{\nu\in\Lambda(\vk',M)}\int_{B_H(R)^c} \mm_\vk(u)\,\nu(\dd u)\le M e^{(\vk-\vk')R^2}\to0\quad\mbox{as $R\to\infty$},
\]
we obtain~\eqref{6.63}. This completes the proof of Theorem~\ref{T:1.1}.
\end{step}


\section{Appendix}\label{S:5}

\subsection{The Foia\c{s}--Prodi estimate}

Let us take any numbers $T, \lambda>0$, any function $\varphi \in L^2([0,T]; H)$, any~integer $N\ge1$, and consider the equations
\begin{align}
\dot u+B(u)+Lu &=h(x)+\p_t\varphi(t,x),\label{E:5.1}\\
\dot v+B(v)+Lv+{\mathsf P}_N[\la(v-u)+B(u)-B(v)]&=h(x)+\p_t\varphi(t,x),\label{E:5.2}
\end{align}
where $\pp$ is the orthogonal projection in~$H$ onto the space $H_N$ defined by~\eqref{ddftgr}. The following result is a version of the Foia\c{s}--Prodi estimate obtained~in~\cite{FP-1967}; see~also~\cite[Section~2.1.8]{KS-book} for a similar result for the Navier--Stokes system (with different equation instead of~\eqref{E:5.2}) and~\cite[Section~7.3]{MN-2015} for the damped nonlinear wave~equation.

\begin{prop}\label{4.13}
Let $u,v\in C([0,T];H) \cap L^2([0,T];H^1) $ be solutions of~\eqref{E:5.1} and~\eqref{E:5.2} issued from $z$ and $z'$, respectively. Then
\begin{equation}\label{FPE1}
\|\pp(u_t-v_t)\| \le e^{- \la t}\|\pp(z-z') \|, \quad t\in [0,T].
\end{equation}
If we assume additionally that
\begin{equation}\label{4.17}
\int_0^t(\| u_s\|_1^2+\| v_s\|_1^2)\,\dd s\le \rho+Kt, \quad t\in [0,T]
\end{equation}
for some numbers $\rho>0$ and $K>0$, then for any $\alpha >0$, we have
\begin{equation}\label{4.16}
\|u_t-v_t \|\leq C_{\la,N} e^{-\alpha t+c \rho}\|z-z'\|, \quad t\in [0,T],
\end{equation}
provided that $2\la> N^2\ge \alpha +cK$. Here $c>0$ is an absolute constant and $C_{\la,N}$ is a constant depending on $\la$ and $N$.
\end{prop}

\begin{proofc}

\begin{step}[Step~1: Proof of~\eqref{FPE1}]
Let us set $y= \pp(u-v)$. Then
\[
\dot y +L y +\la y =0.
\]
Taking the scalar product in $H$ of this equation with $y$, we obtain
\[
\frac12\frac{\dd}{\dd t} \|y\|^2+\|y\|_1^2+\la \| y\|^2=0.
\]
Hence
\[
\frac{\dd}{\dd t} \|y\|^2 +2\la \| y\|^2\le0,
\]
which implies~\eqref{FPE1}.
\end{step}

\begin{proof}[Step~2: Proof of~\eqref{4.16}]
Let $w= u-v$. Then
\begin{equation}\label{lmaeve1}
\dot w+Lw+\la \pp w+ \qq[B(u)-B(v)]=0,
\end{equation}
where $\qq=1-\pp$. For any $a, b \in H^1$, let us set $B(a,b)=\Pi(\langle a,\nabla \rangle b)$. Taking the scalar product of~\eqref{lmaeve1} with $w$, and using the equality
\begin{equation}\label{ezecsd}
B(v)-B(u)=B(v,w)+B(w,u),
\end{equation}
we get
\begin{align}\label{lmaeve}
\frac12\frac{\dd}{\dd t} \|w\|^2+\|w\|_1^2+\la \| \pp w\|^2&= \lag B(v)-B(u),\qq w\rag\nonumber\\&= \lag B(v,w),\qq w\rag+ \lag B(w,u),\qq w\rag\nonumber\\&=:I_1+I_2.
\end{align}
Using the identity
\[
\lag B(a,b),b\rag=0, \quad a,b\in H^1
\]
and the H\"older inequality, we obtain
\begin{multline}\label{lmaeve2}
|I_1|=| \lag B(v, \pp w), \qq w\rag| \le C_1\int_{\T^2} |v| |\nabla \pp w|| \qq w| \dd x \\
\le C_1\|v\| \|\nabla\pp w\|_{\ty} \| w\| \le \frac12 \|\nabla\pp w\|_{\ty}^2+C_2 \|v\|^2 \|w\|^2.
\end{multline}
To estimate $I_2$, we use the H\"older inequality, the inclusion $H^{\frac12}\subset L^4$, and the interpolation inequality $\|a\|_{1/2}^2\le \|a\|\|a\|_1$:
\begin{multline*}
|I_2|=|\lag B(w,u),\qq w\rag|\le C_3\int_{\T^2} |w||\nabla u| |\qq w| \dd x\\
\le C_3 \|w\|_{L^4} \|u\|_1 \|\qq w\|_{L^4} \le C_4 \|w\|\|w\|_1 \|u\|_1\le \frac12 \|w\|_1^2 + C_5 \|w\|^2 \|u\|_1^2.
\end{multline*}
Combining this with~\eqref{lmaeve} and~\eqref{lmaeve2}, and using the Poincar\'e inequality
\[
N\|\qq w\|\le \|w\|_1,
\]
we get
\begin{align}\label{lmaeve3}
\frac{\dd}{\dd t} \|w\|^2+\left(l-c_1(\| u\|_1^2+\| v\|_1^2)\right)\|w\|^2 &\le \|\nabla\pp w\|_{\ty}^2,
\end{align}
where $l=\min\{N^2, 2\la\} $. From~\eqref{FPE1} we deduce that
\[
\|\nabla\pp w_t\|_{\ty}^2\le C_N e^{-2\la t}\|w_0\|^2.
\]
Hence, \eqref{lmaeve3} and~\eqref{4.17} imply that
\[
\|w_t\|^2\le \left(1+ C_N \int_0^t e^{(l-2\la)s} \dd s\right) \|w_0\|^2\exp\left(-l t+c_1 (\rho+Kt)\right).
\]
Choosing $\la$ and $N$ such that $2\la> N^2\ge 2\alpha +c_1K$, we get~\eqref{4.16} with $c=c_1/2$.
\end{proof}
\let\qed\relax
\end{proofc}

\subsection{Proof of Proposition~\ref{P:dist}}\label{apascss}



We closely follow the arguments of the proof of a similar result from Section~7.3 of~\cite{MN-2015} in the case of the nonlinear wave equation (see also~\cite[Section~3.3.3]{KS-book}).


%\smallskip

Note that inequality~\eqref{eoejtnvf} concerns the laws of the solutions and not the solutions themselves. Thus we can choose the underlying probability space~$(\Omega,\FF,\pP)$. We assume that $\Omega$ is the space $C(\R_+; \R)$ endowed with the topology of uniform convergence on bounded intervals, $\pP$ is the law of the Wiener process $W$ in~\eqref{0.2}, and~$\FF$ is the completion of the Borel $\sigma$-algebra of $\Omega$ with respect to $\pP$. We define a stopping~time by
\[
\tau^u=\inf\{t\geq 0:\EE^u(t)\geq \|z\|^2+Kt+\rho\},
\]
where $\EE^u(t)$ is the functional and $K$ is the number in Lemma~\ref{L:5.1}, and $\rho>0$ is a constant to be chosen later. The stopping times $\tau^{u'}$ and $\tau^v$ are defined in a similar way. Then by inequality~\eqref{5.6}, we have
\begin{equation}\label{estd}
\pP\{\tau^{u}<\ty\}+\pP\{\tau^{u'}<\ty\} \le 2e^{-\gamma_0\rho}.
\end{equation}
We define a transformation~$\Lambda:\Omega\to\Omega$ by
\[
\Lambda(\om)(t)=\om(t)-\int_0^t \ph(s,\omega)\dd s,\quad \ph(t,\omega)=\I_{\{t\le\tilde\tau\}}{\mathsf P}_N[\la(v-u)+B(u)-B(v)],
\]
where $
\tilde\tau=\tau^{u}\wedge\tau^{u'}\wedge\tau^{v} $ and $\I_{\{t\leq\tilde\tau\}}$ is the indicator function of the interval $[0,\tilde\tau]$. We use the following result, whose proof is given at the end of this section.

\begin{lemm}\label{4.21}
There is an integer $N_1\ge1$ such that for any numbers $N\ge N_1$, $\la\ge N^2/2 $, and $\rho>0$ and any initial points $z,z'\in H$, we have
\begin{equation}\label{6.46}
\|\Lambda_*\pP-\pP\|_{ var}\leq \left[\exp\left(C_{\lambda,N}\|z-z'\|^2 e^{ C (\|z\|^2 +\|z'\|^2+\rho)} \right)-1\right]^{1/2},
\end{equation}
where $\Lambda_*\pP$ stands for the image of $\pP$ under $\Lambda$, and $C$ and~$C_{\la,N}$ are positive constants not depending on $\rho, z,z'$.
\end{lemm}

Let us introduce auxiliary processes $ y_{u'}$ and $y_v$ in $H$ defined as follows: for~$t\leq \tilde\tau$ they coincide with the processes~$ u'$ and~$v$, respectively, while for $t\geq\tilde\tau$ and $\tilde\tau<\ty$ they are zero. With probability 1, we have
\begin{equation}\label{4.29}
y_v(t,\om)=y_{u'}(t,\Lambda(\om)), \quad t\in J.
\end{equation}
Let us denote by $u_1'$ and $v_1$ the restrictions of $u'(t)$ and $v(t)$ to $J$. Then
\begin{multline*}
\| \nu(z,z')-\nu'(z')\|_{var}=\sup_{\Gamma}|\pP\{v_1\in\Gamma\}-\pP\{u'_1\in\Gamma\}\notag|\\
\leq \pP\{\tilde\tau<\ty\}+ \sup_{\Gamma}|\pP\{v_1\in\Gamma, \tilde\tau=\ty\} -\pP\{u'_1\in\Gamma, \tilde\tau=\ty\}|=\elll_1+\elll_2,
\end{multline*}
where the supremum is taken over all Borel subsets of $C(J;H)$. Note that
\[
\elll_2\leq\|\Lambda_*\pP-\pP\|_{var}.
\]
Further, we have
\[
\elll_1\leq \pP\{\tau^v<\ty, \tau^u\wedge\tau^{u'}=\ty\}+\pP\{\tau^u<\ty\}+\pP\{\tau^{u'}<\ty\}.
\]
Moreover, thanks to~\eqref{4.29},
\begin{align*}
\pP\{\tau^v<\ty, \tau^u\wedge\tau^{u'}=\ty\}&\leq\pP\{\tau^{y_v}<\ty\}=\Lambda_*\pP\{\tau^{y_{u'}}<\ty\}\notag\\
&\leq \pP\{\tau^{y_{u'}}<\ty\}+\|\Lambda_*\pP-\pP\|_{var}\notag\\
&\leq \pP\{\tau^{{u'}}<\ty\}+\|\Lambda_*\pP-\pP\|_{var}.
\end{align*}
Combining last four inequalities, we infer that
\[
\| \nu(z,z')-\nu'(z')\|_{var}\leq 2\left(\pP\{\tau^{u}<\ty\}+\pP\{\tau^{u'}<\ty\}+\|\Lambda_*\pP-\pP\|_{var}\right).
\]
Finally using this with inequalities~\eqref{estd} and~\eqref{6.46}, we get
\[
\| \nu(z,z')-\nu'(z')\|_{var}\leq 4\,e^{-\gamma_0\rho}+2\left[\exp\left(C_{\la, N}\|z-z'\|^2 e^{C(\|z\|^2+\|z'\|^2+\rho)} \right)-1\right]^{1/2}.
\]
Choosing $a=2\gamma_0/(\gamma_0+1)$ and~$\rho=-\gamma_0^{-1}a\ln(\e/4^{1/a})$, we obtain~\eqref{eoejtnvf}.

\goodbreak
\begin{proofc}[Proof of Lemma~\ref{4.21}]

\begin{step}[Step 1: Girsanov theorem]
We write~$ \Omega=\Omega_N\dt+\Omega_N^{\perp}$, where $\Omega_N=C(\R_+; H_N)$ and $\Omega_N^\perp=C(\R_+; { H^\perp_N})$. For any $\om=\om_{1}\dt+\om_{2}\in \Omega$, we write $\om=(\om_{1},\om_{2})\in \Omega_N\times \Omega_N^{\perp}$. Then the transformation $\Lambda$ can be written as $\Lambda(\om)=(\Upsilon(\omega),\om_{2}), $ where $\Upsilon:\Omega\to\Omega_N$ is given by
\[
\Upsilon(\om)(t)=\om_{1}(t)+\int_0^t \ph(s,\om)\dd s.
\]
It is not difficult to see that
\[
\|\Lambda_*\pP-\pP\|_{var}\leq\int_{\Omega^\perp_N}\|\Upsilon_*(\pP_N,\om_{2})-\pP_N\|_{var}\pP^\perp_N(\dd\om_{2}),
\]
where $\pP_N$ and $\pP^\perp_N$ are the images of $\pP$ under the projections $\hat {\mathsf P}_N:\Omega\to\Omega_N$ and $\hat {\mathsf Q}_N:\Omega\to\Omega^\perp_N$, respectively. Let
\[
X=\om_{1}(t),\quad
\hat X=\om_{1}(t)+\int_0^t \ph(s,\om)\dd s.
\]
Then $\pP_N$ coincides with the law $\DD(X)$ of the random variable $X$ and~$\Upsilon_*(\pP_N,\om_{2})$ coincides with that of $\hat X$. By the Girsanov theorem (see~\cite[Theorem~A.10.1]{KS-book}), we have
\begin{equation}\label{Dz-D tilde z}
\|\DD(\hat X)-\DD(X)\|_{var}\leq\frac{1}{2}\left(\left(\E\exp\left[6\max_{1\leq j\leq N} b_j^{-1}\int_0^\infty \|\varphi(t)\|^2\dd t\right]\right)^{\frac{1}{2}}-1\right)^{\frac{1}{2}},
\end{equation}
provided that the Novikov condition
\[
\E \exp\left(p\int_0^\ty \|\ph(t)\|^2 \dd t \right)<\infty \quad \text{ for any $ p>0$}
\]
is satisfied. In Step 2, we show that
\begin{equation}\label{novikov}
\E \exp\left(p\int_0^\ty \|\ph(t)\|^2 \dd t \right)\le \exp\left(C_{p,\lambda, N} \|z-z'\|^2 e^{C(\|z\|^2+\|z'\|^2+\rho)} \right)
\end{equation}
for any $ p>0$. Clearly, this and~\eqref{Dz-D tilde z} imply~\eqref{6.46}.
\end{step}

\begin{proof}[Step 2: Proof of~\eqref{novikov}]
By Proposition~\ref{4.13}, the following inequalities hold
\begin{align}
\|\pp(u_t-v_t)\| &\le e^{- \lambda t}\|\pp(z-z') \|, \quad t\ge0,\label{tildeaz0}
\\
\|u_t-v_t \|&\leq C_{1} e^{- t+c (\|z\|^2+\|z'\|^2+ \rho)}\|z-z'\|, \quad t\in [0,\tilde \tau],\label{tildeaz}
\end{align}
if $2\la>N^2\ge1+cK$. We have
\begin{multline}\label{naxord}
\E\exp\left(p\int_0^\infty \|\ph (t)\|^2\dd t\right) =\E\exp\left(p\int_0^{\tilde\tau} \|\ph(t)\|^2\dd t\right)\\
\leq\E\exp\left(C_2\int_0^{\tilde\tau}\left(\| {\mathsf P}_N[u-v]\|^2+\| {\mathsf P}_N[B(u)-B(v)]\|^2\right)\dd t\right).
\end{multline}
Integrating by parts and using the H\"older inequality, we see that
\[
|\lag B(a,b), e_j\rag| \le C'_j\|a\|_1 \|b\|, \quad a,b\in H^1, \, j\ge1.
\]
Combining this with~\eqref{ezecsd} and~\eqref{tildeaz0}--\eqref{naxord}, we~get
\begin{align*}
\E&\exp\left(p\int_0^\infty \|\ph (t)\|^2\dd t\right)\\
&\leq\E\exp\left(C_3\|z-z'\|^2\int_0^\infty e^{- t+c (\|z\|^2+\|z'\|^2+ \rho)}(1+\|z\|^2+\|z'\|^2+Kt+\rho)^2\dd t\right)\\&\le\E\exp\left(C_4\|z-z'\|^2\int_0^\infty e^{- t/2+2c (\|z\|^2+\|z'\|^2+ \rho)} \dd t\right)\\&=\exp\left(2C_{4}\|z-z'\|^2 e^{C(\|z\|^2+\|z'\|^2+\rho)} \right).
\end{align*}
This proves~\eqref{novikov}.
\end{proof}
\let\qed\relax
\end{proofc}

\subsection{A priori estimates}\label{S:5.1}

The following lemma gathers some standard a priori estimates for the solutions of the stochastic Navier--Stokes system. The reader is referred to Section~2.4.2 in~\cite{KS-book} for more general results.

\begin{lemm}\label{L:5.1}
Assume\,\footnote{Recall that $\BBBB_{i}=\sum_{j\ge1}\alpha_j^i b_j^2$, $i=0,1$.} that $\BBBB_{1}< \ty$, $h\in H^1$, and $u_t$ is a solution of~\eqref{0.4} issued from $u\in H$. Then we have the following estimates.

Exponential moments\pointir There are numbers $\gamma_0=\gamma_0(\BBBB_{0})>0$ and $K=\linebreak K(\BBBB_{0},\|h\|)>0$ such that for any $\varkappa\in(0,\gamma_0)$,
\begin{align}
& \pP_u\left\{\sup_{t\ge0} \left(\EE(t) -Kt\right)\ge \|u\|^2+\rho
\right\} \le e^{-\gamma_0\rho}, \quad \rho\ge0,\label{5.6}
\\
&\E_u e^{\varkappa \EE(t)} \le C_1(\varkappa, \BBBB_0) {e^{\varkappa(Kt+\|u\|^2)}},\label{5.1}
\\
&\E_u \exp(\varkappa \|u_t\|^2)\le e^{-\varkappa t} \exp(\varkappa \|u\|^2)+C_2(\varkappa, \BBBB_{0},\|h\|),\label{5.2}
\\
&\E_u \left\{\|u_t\|^2\exp(\varkappa \|u_t\|^2)\right\}\le C_3(t,\varkappa, \BBBB_{0},\|h\|) \exp(\varkappa \|u\|^2),\label{5.4}
\end{align}
where $\EE(t)=\EE^u(t):= \|u_t\|^2+\int_0^t \|u_s\|_{1}^2\dd s$.

Polynomial moments\pointir For any $m\ge1$,
\begin{align}
&\E_u \|u_t\|^{2m}\le e^{-m\alpha_1 t} \|u\|^{2m}+C_4(m, \BBBB_{0}, \|h\|),\label{5.3}
\\
&\E_u \|u_t\|_1^{2m}\le C_5(t,m, \BBBB_{1},\|h\|_1) \|u\|^{8m}.\label{5.5}
\end{align}
\end{lemm}

\begin{proof}
Estimate~\eqref{5.6} is established in Proposition~2.4.10 in~\cite{KS-book}. To\linebreak prove~\eqref{5.1}, we denote $\tilde\CC_\rho:=\CC_{\rho-1}\setminus\CC_\rho $, where $\CC_\rho$ the event on the left-hand side of~\eqref{5.6} and $\CC_{-1}:=\Omega$. Then for any $\vk\in (0,\gamma_0)$, we have
\begin{align*}
\E_u e^{\varkappa \EE(t)} &= \sum_{\rho=0}^\ty\E_u \left\{e^{\varkappa \EE(t)} I_{\tilde \CC_\rho}\right\}\le e^{\varkappa(Kt+\|u\|^2)}
\sum_{\rho=0}^\ty e^{\varkappa \rho} \pP\{ \CC_{\rho-1}\} \\&\le e^{\varkappa(Kt+\|u\|^2)+\gamma_0}
\sum_{\rho=0}^\ty e^{(\varkappa-\gamma_0) \rho} = \frac{e^{\varkappa(Kt+\|u\|^2)+\gamma_0}}{ 1 -e^{(\varkappa-\gamma_0) }}.
\end{align*}
Estimates~\eqref{5.2} and~\eqref{5.5} are proved in Propositions~2.4.9 and~2.4.12 in~\cite{KS-book}, respectively. To show\,\footnote{We confine ourselves to a formal derivation of~\eqref{5.3}. The accurate proof is based on the same arguments applied to the stopped solutions $u(t\wedge \tau_n)$, where $\tau_n=\inf\{t\ge0: \|u(t)\|> n\}$.}~\eqref{5.3}, we set~$F(u)= \|u\|^{2m}$. Then
\begin{align*}
\p_u F(u;v)&=2m \|u\|^{2(m-1)} \lag u,v\rag,\\
\p^2_u F(u;v)&= 2m\|u\|^{2(m-1)} \|v\|^2+4m(m-1)\|u\|^{2(m-2)} \lag u,v\rag^2,
\end{align*}
so applying the It\^o formula for the functional~$F$ and taking the expectation:
\begin{multline*}
\E_u\|u_t\|^{2m}=\|u\|^{2m}+\E_u\int_0^t \Bigg(2m\|u_s\|^{2(m-1)}\lag u_s, -L u_s-B(u_s)+h\rag\\
+ m \|u_s\|^{2(m-1)} \BBBB_{0}+2m(m-1)\|u_s\|^{2(m-2)} \sum_{j=1}^\ty b_j^2u_j^2
\Bigg)\dd s,
\end{multline*}
where $u_j=\lag u, e_j\rag$. The identity
\begin{equation}\label{5.7}
\lag u, B(u)\rag=0
\end{equation}
and the Cauchy--Schwarz and Poincar\'e inequalities imply that
\begin{align*}
\E_u\|u_t\|^{2m}&\le\|u\|^{2m}+\E_u\int_0^t \big(2m\|u_s\|^{2(m-1)}(-\|u_s\|_1^2+\| u_s\|\|h\|) \\
&\qquad\qquad\qquad+ m \|u_s\|^{2(m-1)} \BBBB_{0}+2m(m-1)\|u_s\|^{2(m-1)} \BBBB_{0}\big)\dd s\\
&\le\|u\|^{2m}-m\alpha_1 \int_0^t \E_u\|u_s\|^{2m}\dd s +tC_6(m,\BBBB_{0},\|h\|).
\end{align*}
Combining this with the Gronwall inequality, we obtain~\eqref{5.3}.

%\smallskip
To prove~\eqref{5.4}, we apply the It\^o formula for $F(t,u)=t\|u\|^2\exp(\varkappa \|u\|^2)$, use the equalities
\begin{align*}
\p_t F(t,u;v)&=\|u\|^2\exp(\varkappa \|u\|^2),\\
\p_u F(t,u;v)&=2t\exp(\varkappa \|u\|^2)(1+ \varkappa\|u\|^2)\lag u,v\rag,\\
\p^2_u F(t,u;v)&=2t\exp(\varkappa \|u\|^2)\big(2\vk (2+ \varkappa\|u\|^2)\lag u,v\rag ^2 + (1+ \varkappa\|u\|^2)\|v\|^2 \big),
\end{align*}
and take the expectation:
\begin{multline*}
t\E_u\left\{\|u_t\|^2 \exp(\varkappa \|u_t\|^2)\right\}= \E_u\!\int_0^t\!\! \Bigg(\|u_s\|^2\!+\!2s(1\!+\! \varkappa\|u_s\|^2)\lag u_s, \!-L u_s\!-\!B(u_s)\!+\!h\rag \\
+s\sum_{j=1}^\ty\Big[2\vk (2+ \varkappa\|u_s\|^2)b_j^2u_j^2 + (1+ \varkappa\|u_s\|^2)b_j^2 \Big]\Bigg) \exp(\varkappa \|u_s\|^2) \dd s.
\end{multline*}
Again using~\eqref{5.7} and the Cauchy--Schwarz and Poincar\'e inequalities, we get for sufficiently small $\gamma_0=\gamma_0(\BBBB_{0})>0$ and any~$\vk\in (0,\gamma_0)$,
\begin{multline*}
t\E_u\left\{\|u_t\|^2\exp(\varkappa \|u_t\|^2)\right\}\le \E_u\int_0^t \big(\|u_s\|^2+ 2s(1+ \varkappa\|u_s\|^2) (-\|u_s\|^2_1+ \|u_s\|\|h\|) \\
+ s 2\vk (2+ \varkappa\|u_s\|^2)\|u_s\|^2\BBBB_{0} + s(1+ \varkappa\|u_s\|^2)\BBBB_{0} \big)\exp(\varkappa \|u_s\|^2) \dd s\\
\le \E_u\int_0^t \Big(\|u_s\|^2 +sC_7(\varkappa, \BBBB_{0}, \|h\|) \Big)\exp(\varkappa \|u_s\|^2) \dd s.
\end{multline*}
Thus~\eqref{5.4} follows from~\eqref{5.2}, the Poincar\'e inequality, and the estimate
\[
\E_u \left\{\int_0^t\|u_s\|_1^2\exp(\varkappa \|u_s\|^2)\dd s\right\} \le C_8(t,\varkappa, \BBBB_{0},\|h\|) \exp(\varkappa \|u\|^2).
\]
The latter is easily proved by applying the It\^o formula for $F(u)=\exp(\varkappa \|u\|^2)$. This completes the proof of the lemma.
\end{proof}

\subsection{Hyper-exponential recurrence}

For any $R>0$, let~$\tau(R)$ be the first hitting time of the set~$X_R$:
\begin{equation}\label{E:tau}
\tau(R)= \inf\{t \ge0:u_{t}\in X_R\}.
\end{equation}
We have the following standard estimate for the exponential moment of $\tau(R)$.

\begin{prop}\label{P:2.1}
For any $\gamma>0$, there are positive numbers $m,R,$ and $C$ such~that
\begin{equation}\label{8.1}
\E_u\exp\bigl(\gamma \tau(R)\bigr)\le C\,\wwww_m(u),
\quad u\in H.
\end{equation}
\end{prop}

\begin{proof}
See~\cite[Proposition~5.1]{JNPS-2014} for a similar result in the discrete-time case. The proof of~\eqref{8.1} follows the same arguments. The idea is to establish the inequality for the first hitting time of a ball in~$H$ and then to use the regularising property of the Navier--Stokes system.

\begin{step}[Step~1: Hyper-exponential recurrence in $H$]
For any $r>0$, we denote by $\tau_0(r)$ the first hitting time of the ball~$B_{H}(r)$:
\[
\tau_0(r)=\inf\{t\ge0:u_{t}\in B_{H}(r)\}.
\]
Let us prove that, for any $\gamma>0$, there are positive numbers $m,r,$ and $C$ such~that
\begin{equation}\label{8.1a}
\E_u\exp\bigl(\gamma \tau_{0}(r)\bigr)\le C\,\wwww_m(u),
\quad u\in H.
\end{equation}
Indeed, let $m>0$ be so large that $q:=2e^{-m\alpha_1 }<1$. Then, by~\eqref{5.3}, we have
\begin{equation}\label{8.5}
\E_u\|u_1\|^{2m}\le q\,\bigl(\|u\|^{2m}\vee r\bigr),\quad u\in H,
\end{equation}
where $r=e^{m\alpha_1 } C_4(m, B_0, \|h\|).$ The Markov property and~\eqref{8.5} imply that
\[
p_k(u):=\E_u\bigl(I_{\{\tau_0(r)>k\}}\|u_k\|^{2m}\bigr)\le q^k\|u\|^{2m},\quad k\ge0,\, u\in H
\]
(cf.~\cite[proof of Lemma~3.6.1]{KS-book}),
hence
\begin{equation}\label{8.6}
\IP_u\{\tau_0(r)>k\}\le r^{-2m}p_k(u)\le r^{-2m} q^k\|u\|^{2m}.
\end{equation}
As $q<1$, the Borel--Cantelli lemma gives that $\IP_u\{\tau_0(r)<\ty\}=1$. Choosing~$m\ge1$ so large that $e^{\gamma} q<1$, we derive from~\eqref{8.6}
\begin{align*}
\E_u\exp\bigl(\gamma \tau_0(r)\bigr)&\le1+\sum_{k=1}^\infty e^{\gamma k}\IP_u\{\tau_0(r)>k-1\}\\
&\le1+r^{-2m}\|u\|^{2m}\sum_{k=1}^\infty e^{\gamma k}q^{k-1}\le C\,\wwww_m(u),
\end{align*}
which proves~\eqref{8.1a}.
\end{step}

%\smallskip
\begin{proof}[Step~2: Hyper-exponential recurrence in $H^1$]
First note that, for any numbers $p\in (0,1)$ and $ r>0$, there is~$R>0$ such that
\begin{equation}\label{Eq:ddfg}
\IP_u\bigl\{u_1\in X_R\bigr\}\ge 1-p, \quad u\in B_H(r).
\end{equation}
Indeed, this follows immediately from the Chebyshev inequality and~\eqref{5.5}:
\[
\IP_u\bigl\{\|u_1\|_1>R\big\}\le R^{-2} \E_u \|u_1\|_1^2\le C R^{-2} \|u\|^8\le C R^{-2} r^8\le p
\]
for any $u\in B_H(r)$ and sufficiently large $R=R(r,p)$.

%\smallskip
Now we combine~\eqref{8.1a} and~\eqref{Eq:ddfg} to prove~\eqref{8.1}. We introduce the sequences of stopping times
\[
\tau_0'=\tau_0(r), \quad \tau_n'=\inf\{t\ge \tau_{n-1}'+1: u_t\in B_H(r)\}, \quad n\ge1
\]
and $\tau_n=\tau'_n+1$. Let
\[
\hat n=\min\{n\ge0: \tau_n\in X_R\}.
\]
From~\eqref{Eq:ddfg} and the strong Markov property we have
\[
\IP_u\{\hat n > k\} \le (1 - p)^k, \quad k \ge 0,\, u \in H,
\]
so $\hat n$ is almost surely finite. For any integers $k,M\ge1$, we can write
\begin{align*}
\IP_u\bigl\{\tau(R)\ge M\bigr\}&=\IP_u\bigl\{\tau(R)\ge M,\tau_k<M\bigr\}+\IP_u\bigl\{\tau(R)\ge M,\tau_k\ge M\bigr\}
\notag\\
&\le \IP_u\bigl\{\tau(R)>\tau_k\bigr\}+\IP_u\bigl\{\tau_k\ge M\bigr\}.
\end{align*}
Since $\{\tau(R)>\tau_k\}\subset\{\hat n>k\}$, the first probability is estimated by $(1 - p)^k$. The second one is estimated using~\eqref{8.1a} and the strong Markov property
\[
\IP_u\bigl\{\tau_k\ge M\bigr\}\le C_1^k\,\wwww_m(u)e^{-3\gamma M},
\]
where $C_1>0$ does not depend on $k, M\ge1$ and $u\in H$.
Thus, we obtain
\[
\IP_u\bigl\{\tau(R)\ge M\bigr\}\le (1-p)^k+C_1^k\,\wwww_m(u)e^{-3\gamma M}.
\]
To complete the proof, it remains to choose appropriately the parameters $k$ and~$R$.
We take $k\sim \varepsilon M$, where $\varepsilon>0$ is so small that $\varepsilon\log C_1\le \gamma$, and $R>0$ so large that~$\varepsilon\log(1-p)^{-1}\ge2\gamma$. Then
\[
\IP_u\bigl\{\tau(R)\ge M\bigr\}\le 2e^{-2\gamma M}\wwww_m(u),
\]
which implies~\eqref{8.1}.
\end{proof}
\let\qed\relax
\end{proof}

\subsection{Generalised Markov semigroups}\label{S:5.2}

For the reader's convenience, we recall here a result on the large-time asymptotics of generalised Markov semigroups in a Polish space $X$. It is established in~\cite{JNPS-2014} in the discrete-time setting, then extended to the continuous-time in~\cite{MN-2015}. Let us first recall some terminology.

\begin{defi}\label{D:5.2}
We shall say that $\{P_t(u,{\cdot\,}),u\in X, t\ge0\} $ is a \emph{generalised Markov family of transition kernels} if the following two properties are satisfied.

Feller property\pointir For any $t\ge0$, the function $u\mapsto P_t(u,{\cdot\,})$ is continuous from~$X$ to~$\MM_+(X)$ and does not vanish.

Kolmogorov--Chapman relation\pointir For any $t,s\ge0, u\in X$, and Borel set~$\Gamma\subset X$, the following relation holds
\[
P_{t+s}(u,\Gamma)=\int_X P_s(v,\Gamma) P_t(u,\dd v).
\]
\end{defi}

To any such family we associate two semigroups by the following relations:
\begin{align*}
\PPPP_t&:C_b(X)\to C_b(X), & \PPPP_t \psi(u)&=\int_X \psi(v) P_t(u,\dd v),\\
\PPPP_t^*&:\MM_+(X)\to \MM_+(X), & \PPPP_t^* \mu(\Gamma)&=\int_X P_t(v,\Gamma) \mu(\dd v),\quad t\ge0.
\end{align*}
For a measurable function ${\wwww}:X\to[1,+\infty]$ and a family $\CC\subset C_b(X)$, we denote by~$\CC^\wwww$ the set of functions $\psi\in L_\wwww^\infty(X)$ that can be approximated with respect to the norm~$\|{\,\cdot\,}\|_{L_\wwww^\infty}$ by finite linear combinations of functions from~$\CC$. We shall say that a family $\CC\subset C_b(X)$ is \emph{determining} if for any~$\mu,\nu\in\MM_+(X)$ satisfying $\lag \psi,\mu\rag=\lag \psi,\nu\rag$ for all~$\psi\in\CC$, we have $\mu=\nu$. Finally, a family of functions~$\psi_t: X\to\R$ is \emph{uniformly equicontinuous} on a subset~$K\subset X$ if for any~$\varepsilon>0$ there is~$\delta>0$ such that~$|\psi_t(u)-\psi_t(v)|<\varepsilon$ for any~$u\in K$, $v\in B_X(u,\delta)\cap K$, and~$ t\ge0$. The following result is Theorem~7.4 in~\cite{MN-2015}.

\begin{theo}\label{T:5.3}
Let~$\{P_t(u,{\cdot\,}),u\in X, t\ge0\} $ be a generalised Markov family of transition kernels satisfying the following four properties.

Growth conditions\pointir There is an increasing sequence~$\{X_R\}_{R=1}^\infty$ of compact~subsets of~$X$ such that~$X_\infty:=\bigcup_{R=1}^\ty X_R$ is dense in~$X$. The measures~$P_{t}(u,{\cdot\,})$ are concentrated on~$X_\infty$ for any $u\in X$ and $t>0$, and there is a measurable function $\wwww:X\to[1,+\infty]$ and an integer~$R_0\ge1$ such that\,\footnote{The expression $(\PPPP_t\wwww)(u)$ is understood as an integral of a positive function~$\wwww$ against a positive measure $P_t(u,{\cdot\,})$.}
\begin{align*}
&\sup_{t\ge0}
\frac{\|\PPPP_t\wwww\|_{L_\wwww^\infty}}{\|\PPPP_t{\mek}\|_{R_0}}<\infty,\\
&\sup_{t\in [0,1]} \|\PPPP_t{\mek}\|_{\infty}<\ty,
\end{align*}
where $\|{\,\cdot\,}\|_R$ and $\|{\,\cdot\,}\|_\ty$ denote the~$L^\infty$ norm on~$X_R$ and $X$, respectively, and we set $\infty/\infty=0$.

Time-continuity\pointir For any $g\in C_\wwww(X)$ and $u\in X$, the function~$t\mapsto\PPPP_tg(u) $ is continuous from~$\R_+$ to $\R$.

Uniform irreducibility\pointir For sufficiently large $\rho\ge1$, any $R\ge 1$ and $r>0$, there are positive numbers $l=l(\rho,r,R)$ and~$p=p(\rho,r)$ such that
\[
P_l(u,B_{X}(\hat u,r))\ge p\quad\mbox{for all~$u\in X_R,\,\hat u\in X_\rho$}.
\]

Uniform Feller property\pointir There is a number~$R_0\ge1$ and a determining family~$\CC\subset C_b(X)$ such that~$\mek\in\CC$ and the family $\{\|\PPPP_t{\mek}\|_R^{-1}\PPPP_t\psi, t\ge0\}$ is uniformly equicontinuous on~$X_R$ for any~$\psi\in \CC$ and~$R\ge R_0$.

Then~for any $t>0$, there is at most one measure $\mu_t\in{\mathcal P}_\wwww(X)$ such that $\mu_t(X_\ty)=1$ and
\[
\PPPP^*_t \mu_t=\la(t)\mu_t \quad \text{for some $\la(t)\in \R$}
\]
satisfying the following condition:
\[
\|\PPPP_t\wwww\|_R\int_{X\setminus X_R}\wwww(u)\, \mu_t(\dd u)\to0
\quad\mbox{as $R\to\infty$}.
\]
Moreover, if such a measure~$\mu_t$ exists for all $t>0$, then it is independent of~$t$~(we set~$\mu:=\mu_t$), the corresponding eigenvalue is of the form~$\lambda(t)=\la^t$, $\la>0$, $\supp \mu=X$, there is a non-negative function~$h\in L^\ty_\wwww(X)$ such that~$\lag h,\mu\rag=1$,
\[
(\PPPP_t h)(u)=\lambda^t h(u)\quad\mbox{for $u\in X$, $t>0$},
\]
the restriction of~$h$ to~$X_R$ belongs to $C_+(X_R)$,
and for any $\psi\in \CC^\wwww$ and $R\ge1$, we have
\begin{equation}
\lambda^{-t}\PPPP_t \psi\to\lag \psi,\mu\rag h
\quad\text{in~$C(X_R)\cap L^1(X,\mu)$ as~$t\to\infty$}. \label{5.12}
\end{equation}
Finally, if a Borel set $B\subset X$ is such that
\begin{equation}
\sup_{u\in B} \int_{X\setminus X_R}\wwww(v)\,P_s(u,\dd v) \to0
\quad\text{as $R\to\infty$} \label{5.13}
\end{equation}
for some $s>0$,
then
\[
\lambda^{-t}\PPPP_t \psi\to\lag \psi,\mu\rag h
\quad\text{in~$L^\ty(B)$ as~$t\to\infty$}
\]
for any $\psi\in \CC^\wwww$.
\end{theo}


\begin{thebibliography}{99}
\bibitem{BKL-2002}Bricmont, Jean; Kupiainen, Antti; Lefevere, Raphaël Exponential mixing of the 2D stochastic Navier–Stokes dynamics, Commun. Math. Phys., Volume 230 (2002) no. 1, pp. 87-132 | MR | DOI | Zbl
\bibitem{DS1989}Deuschel, Jean-Dominique; Stroock, Daniel W. Large Deviations, Pure and Applied Mathematics, 137, Academic Press Inc., 1989 | Zbl
\bibitem{dudley2002}Dudley, Richard M. Real analysis and probability, Cambridge Studies in Advanced Mathematics, 74, Cambridge University Press, 2002 | MR | Zbl
\bibitem{DV-1975}Donsker, Monroe D.; Varadhan, S. R. Srinivasa Asymptotic evaluation of certain Markov process expectations for large time. I–II, Commun. Pure Appl. Math., Volume 28 (1975) no. 1, p. 1-47 \& 279–301 | DOI | MR | Zbl
\bibitem{ADOZ00}Dembo, Amir; Zeitouni, Ofer Large deviations Techniques and applications, Springer, 2000
\bibitem{EMS-2001}E, Weinan; Mattingly, Jonathan C.; Sinaĭ, Yakov G. Gibbsian dynamics and ergodicity for the stochastically forced Navier–Stokes equation, Commun. Math. Phys., Volume 224 (2001) no. 1, pp. 83-106 | MR | Zbl
\bibitem{FM-1995}Flandoli, Franco; Maslowski, Bohdan Ergodicity of the 2D Navier–Stokes equation under random perturbations, Commun. Math. Phys., Volume 172 (1995) no. 1, pp. 119-141 | DOI | Zbl
\bibitem{FP-1967}Foiaš, Ciprian; Prodi, Giovanni Sur le comportement global des solutions non-stationnaires des équations de Navier–Stokes en dimension 2, Rend. Semin. Mat. Univ. Padova, Volume 39 (1967), pp. 1-34 | Numdam | Zbl
\bibitem{FW1984}Freidlin, Mark I.; Wentzell, Alexander D. Random perturbations of dynamical systems, Grundlehren der Mathematischen Wissenschaften, 260, Springer, 1984 | MR | Zbl
\bibitem{gourcy-2007b}Gourcy, Mathieu A large deviation principle for 2D stochastic Navier–Stokes equation, Stochastic Processes Appl., Volume 117 (2007) no. 7, pp. 904-927 | DOI | MR | Zbl
\bibitem{gourcy-2007a}Gourcy, Mathieu Large deviation principle of occupation measure for a stochastic Burgers equation, Ann. Inst. Henri Poincaré, Probab. Stat., Volume 43 (2007) no. 4, pp. 375-408 | Numdam | Zbl | MR
\bibitem{HM-2006}Hairer, Martin; Mattingly, Jonathan C. Ergodicity of the 2D Navier–Stokes equations with degenerate stochastic forcing, Ann. Math., Volume 164 (2006) no. 3, pp. 993-1032 | MR | DOI | Zbl
\bibitem{JNPS-2012}Jakšić, Vojkan; Nersesyan, Vahagn; Pillet, Claude-Alain; Shirikyan, Armen Large deviations from a stationary measure for a class of dissipative PDEs with random kicks, Commun. Pure Appl. Math., Volume 68 (2015) no. 12, pp. 2108-2143 | MR | DOI | Zbl
\bibitem{JNPS-2014}Jakšić, Vojkan; Nersesyan, Vahagn; Pillet, Claude-Alain; Shirikyan, Armen Large deviations and mixing for dissipative PDE’s with unbounded random kicks, Nonlinearity, Volume 31 (2018) no. 2, pp. 540-596 | MR | DOI | Zbl
\bibitem{KNS-2018}Kuksin, Sergei; Nersesyan, Vahagn; Shirikyan, Armen Exponential mixing for a class of dissipative PDEs with bounded degenerate noise (2018) https://arxiv.org/abs/1802.03250 | Zbl
\bibitem{KS1991}Karatzas, Ioannis; Shreve, Steven E. Brownian motion and stochastic calculus, Graduate Texts in Mathematics, 113, Springer, 1991 | MR | Zbl
\bibitem{KS-cmp2000}Kuksin, Sergei; Shirikyan, Armen Stochastic dissipative PDEs and Gibbs measures, Commun. Math. Phys., Volume 213 (2000) no. 2, pp. 291-330 | MR | DOI | Zbl
\bibitem{KS-jmpa2002}Kuksin, Sergei; Shirikyan, Armen Coupling approach to white-forced nonlinear PDEs, J. Math. Pures Appl., Volume 81 (2002) no. 6, pp. 567-602 | MR | DOI | Zbl
\bibitem{KS-book}Kuksin, Sergei; Shirikyan, Armen Mathematics of two-dimensional turbulence, Cambridge Tracts in Mathematics, 194, Cambridge University Press, 2012 | MR | Zbl
\bibitem{kuksin-jmpa}Kuksin, Sergei B. Ergodic theorems for 2D statistical hydrodynamics, Rev. Math. Phys., Volume 14 (2002) no. 6, pp. 585-600 | MR | DOI | Zbl
\bibitem{lions1969}Lions, Jacques-Louis Quelques méthodes de résolution des problèmes aux limites non linéaires, Études mathématiques, Dunod, 1969 | Zbl
\bibitem{MN-2015}Martirosyan, Davit; Nersesyan, Vahagn Local large deviations principle for occupation measures of the stochastic damped nonlinear wave equation, Ann. Inst. Henri Poincaré, Probab. Stat., Volume 54 (2018) no. 4, pp. 2002-2041 | MR | DOI | Zbl
\bibitem{MN-2017}Martirosyan, Davit; Nersesyan, Vahagn Multiplicative ergodic theorem for a non-irreducible random dynamical system (2018) https://www.archives-ouvertes.fr/hal-01695046v1
\bibitem{odasso-2008}Odasso, Cyril Exponential mixing for stochastic PDEs: the non-additive case, Probab. Theory Relat. Fields, Volume 140 (2008) no. 1-2, pp. 41-82 | MR | DOI | Zbl
\bibitem{shirikyan-ptrf}Shirikyan, Armen Law of large numbers and central limit theorem for randomly forced PDE’s, Probab. Theory Relat. Fields, Volume 134 (2006) no. 2, pp. 215-247 | MR | DOI | Zbl
\bibitem{wu-2001}Wu, Liming Large and moderate deviations and exponential convergence for stochastic damping Hamiltonian systems, Stochastic Processes Appl., Volume 91 (2001) no. 2, pp. 205-238 | MR | Zbl
\end{thebibliography}
\end{document}
